\chapter{Characteristic Classes}\label{characteristic-classes}
\chapterstatus{complete}

\section{Introduction}\label{characteristic-classes:section-introduction}

Characteristic classes measure how far a vector bundle is from being trivial. They are cohomology classes $c(E) \in H^\bullet(M)$ attached to vector bundles, natural
under pullback ($c(f^*E) = f^*c(E)$) and trivial on trivial bundles; so a bundle with a nonzero characteristic class is not trivial, and in many cases the classes detect
more precise obstructions, such as the impossibility of finding a nowhere vanishing section. The first example is the first Chern class of a complex line bundle
(Theorem~\ref{vector-bundles:theorem-line-bundles}), which classifies line bundles completely.

There are two ways to construct them. The topological way uses the Thom isomorphism or classifying spaces. The differential-geometric way, \emph{Chern--Weil theory}, starts
from a connection: its curvature $F$ is a matrix of $2$-forms, and applying to it a polynomial that is invariant under conjugation gives a closed form whose cohomology class
does not depend on the connection. For the tautological line bundle of $\CC\PP^1$ this was already visible in Chapter~\ref{connections}, where the curvature of a natural
connection integrates to $\int\frac{i}{2\pi}F = -1$, a number that cannot change when the connection changes; and for the tangent bundle of a surface the same mechanism
is the Gauss--Bonnet theorem, $\int K\,dA = 2\pi\chi$.

Characteristic classes are central to the Hodge conjecture: Chern classes of algebraic vector bundles are Hodge classes, the Chern character appears in the Riemann--Roch
theorem, and the class of a divisor is the first Chern class of a line bundle (Chapter~\ref{divisors}).

The chapter goes as follows.
\begin{enumerate}
\item Chern--Weil theory: invariant polynomials, the closedness and independence of the Chern--Weil forms with complete proofs, Chern forms and Chern classes, computed
for line bundles on $\CC\PP^1$ and for the tangent bundle of the sphere (Section~\ref{characteristic-classes:section-chern-weil}).
\item The axioms for Chern classes, the splitting principle, the Chern character and the Todd class (Section~\ref{characteristic-classes:section-axioms}).
\item The Thom isomorphism and the Euler class in de Rham cohomology (Section~\ref{characteristic-classes:section-thom}).
\item Euler, Pontryagin and Stiefel--Whitney classes (Section~\ref{characteristic-classes:section-euler}).
\item The Euler form of a metric connection and Chern's intrinsic proof of the Gauss--Bonnet--Chern theorem: the integral of the Euler form counts the zeros of a section, with
signs (Section~\ref{characteristic-classes:section-gauss-bonnet-chern}).
\item Wu's formula (Section~\ref{characteristic-classes:section-wu}).
\end{enumerate}
References: \cite{milnor-stasheff}, \cite{bott-tu}, \cite{hirzebruch-topological}, \cite[Chapter 1]{griffiths-harris}.

\noindent\textbf{Prerequisites.} Chapter~\ref{connections} (Connections and Curvature), Chapter~\ref{de-rham} (de Rham Cohomology)
and Chapter~\ref{cohomology-operations} (Cohomology Operations).

\section{Chern--Weil theory}\label{characteristic-classes:section-chern-weil}

\begin{definition}\label{characteristic-classes:definition-invariant-polynomial}
An \emph{invariant polynomial} of degree $k$ is a symmetric $k$-linear form $\tilde{P}$ on $\mathfrak{gl}_r(\CC) = M_r(\CC)$ with
$\tilde{P}(gA_1g^{-1}, \ldots, gA_kg^{-1}) = \tilde{P}(A_1, \ldots, A_k)$ for $g \in \GL_r(\CC)$; we write $P(A) = \tilde{P}(A, \ldots, A)$. Differentiating the invariance along $g = \exp(tB)$ gives
\[
\sum_{i=1}^k\tilde{P}(A_1, \ldots, [B, A_i], \ldots, A_k) = 0 .
\]
Multilinear forms are extended to matrices of differential forms by multiplying the entries with the wedge product; then $\tilde{P}(\alpha_1, \ldots, \alpha_k)$ is a form of degree $\sum\deg\alpha_i$, and it is
symmetric in the graded sense.
\end{definition}

\begin{example}\label{characteristic-classes:example-invariant-polynomials}
\begin{enumerate}
\item The trace, $P(A) = \operatorname{tr}A$, of degree $1$, and more generally $P_k(A) = \operatorname{tr}(A^k)$, with polarisation
$\tilde P_k(A_1, \ldots, A_k) = \frac1{k!}\sum_{\sigma \in S_k}\operatorname{tr}(A_{\sigma(1)} \cdots A_{\sigma(k)})$; invariance follows from $\operatorname{tr}(gBg^{-1}) = \operatorname{tr}B$.
\item The determinant, of degree $r$, and more generally the coefficients $\sigma_k(A)$ of the characteristic polynomial,
$\det(I + tA) = \sum_{k=0}^r\sigma_k(A)t^k$; they are invariant because $\det(I + tgAg^{-1}) = \det(g(I + tA)g^{-1}) = \det(I + tA)$. In terms of the eigenvalues $\lambda_j$ of $A$,
$\sigma_k(A)$ is the $k$-th elementary symmetric function and $\operatorname{tr}(A^k) = \sum_j\lambda_j^k$, the $k$-th power sum. So $\sigma_1 = \operatorname{tr}$,
$\sigma_2(A) = \frac12((\operatorname{tr}A)^2 - \operatorname{tr}(A^2))$, and every invariant polynomial is a polynomial in the $\sigma_k$ (a conjugation invariant function of a
diagonalisable matrix is a symmetric function of its eigenvalues, and diagonalisable matrices are dense).
\item On the Lie algebra $\mathfrak{so}(2m)$ of skew-symmetric matrices there is one more invariant for $\mathrm{SO}(2m)$ (not for $\GL$, nor for $\mathrm{O}(2m)$), the
\emph{Pfaffian} (Definition~\ref{multilinear-algebra:definition-pfaffian}), with $\mathrm{Pf}(A)^2 = \det A$ and $\mathrm{Pf}(g^TAg) = \det(g)\mathrm{Pf}(A)$; for $m = 1$,
$\mathrm{Pf}\bigl(\begin{smallmatrix} 0 & a\\ -a & 0\end{smallmatrix}\bigr) = a$. It gives the Euler form of Section~\ref{characteristic-classes:section-gauss-bonnet-chern}.
\end{enumerate}
\end{example}

The key computation in Chern--Weil theory is that invariant polynomials, applied to matrices of forms, turn covariant derivatives into ordinary ones.

\begin{lemma}\label{characteristic-classes:lemma-invariant-derivative}
Let $\tilde P$ be an invariant symmetric $k$-linear form, $\theta$ an $r \times r$ matrix of $1$-forms on an open set, and $\alpha_1, \ldots, \alpha_k$ matrices of forms of degrees
$d_1, \ldots, d_k$. Put $D\alpha = d\alpha + \theta \wedge \alpha - (-1)^{\deg\alpha}\alpha \wedge \theta$. Then
\[
d\tilde P(\alpha_1, \ldots, \alpha_k) = \sum_{i=1}^k(-1)^{d_1 + \cdots + d_{i-1}}\tilde P(\alpha_1, \ldots, D\alpha_i, \ldots, \alpha_k).
\]
\end{lemma}

\begin{proof}
Since $\tilde P$ is multilinear in the matrix entries with constant coefficients, and the entries are multiplied with the wedge product, the Leibniz rule for $d$ gives
$d\tilde P(\alpha_1, \ldots, \alpha_k) = \sum_i(-1)^{d_1 + \cdots + d_{i-1}}\tilde P(\alpha_1, \ldots, d\alpha_i, \ldots, \alpha_k)$. It remains to show that the terms with $\theta$ add up to zero.
Write $\theta = \sum_jdx^j\,\theta_j$ with matrices of functions $\theta_j$. Then $\theta \wedge \alpha_i = \sum_jdx^j \wedge \theta_j\alpha_i$ and
$(-1)^{d_i}\alpha_i \wedge \theta = \sum_j(-1)^{d_i}(-1)^{d_i}dx^j \wedge \alpha_i\theta_j = \sum_jdx^j \wedge \alpha_i\theta_j$, so the extra term in $D\alpha_i$ is $\sum_jdx^j \wedge [\theta_j, \alpha_i]$. Moving
$dx^j$ in front of $\alpha_1, \ldots, \alpha_{i-1}$ produces exactly the sign $(-1)^{d_1 + \cdots + d_{i-1}}$, so the sum of the extra terms is
$\sum_jdx^j \wedge \sum_i\tilde P(\alpha_1, \ldots, [\theta_j, \alpha_i], \ldots, \alpha_k)$. Each inner sum vanishes: writing each $\alpha_i$ as a sum of products of matrices of functions with
scalar forms and using multilinearity, it reduces to the infinitesimal invariance $\sum_i\tilde P(A_1, \ldots, [B, A_i], \ldots, A_k) = 0$ of
Definition~\ref{characteristic-classes:definition-invariant-polynomial}, since the scalar forms can be moved out of $\tilde P$ (with the same signs in every term).
\end{proof}

\begin{theorem}[Chern--Weil]\label{characteristic-classes:theorem-chern-weil}
Let $E \to M$ be a complex vector bundle of rank $r$, $\nabla$ a connection with curvature $F$, and $P$ an invariant polynomial of degree $k$. Then:
\begin{enumerate}
\item $P(F)$ is a well-defined global $2k$-form (computed in any local frame) and $dP(F) = 0$;
\item the class $[P(F)] \in H^{2k}_{\mathrm{dR}}(M)$ does not depend on $\nabla$;
\item the resulting map from invariant polynomials to $H^\bullet_{\mathrm{dR}}(M)$ is natural: $[P(F_{f^*\nabla})] = f^*[P(F_\nabla)]$ for smooth $f \colon N \to M$.
\end{enumerate}
\end{theorem}

\begin{proof}
(1) Under a change of frame the curvature matrix transforms by conjugation, $F^{\prime} = g^{-1}Fg$ (Proposition~\ref{connections:proposition-curvature}), and $P$ is invariant under
conjugation (for matrices of $2$-forms, whose entries commute, the invariance of $P$ on matrices of numbers implies it pointwise); so $P(F)$ is the same in all frames and defines a
global form. By Lemma~\ref{characteristic-classes:lemma-invariant-derivative} with all $\alpha_i = F$, $dP(F) = \sum_i\tilde P(F, \ldots, DF, \ldots, F)$, and $DF = dF + \theta \wedge F - F \wedge \theta = 0$
by the Bianchi identity (Proposition~\ref{connections:proposition-curvature}). So $P(F)$ is closed.

(2) Let $\nabla_0$, $\nabla_1$ be connections and $\alpha = \nabla_1 - \nabla_0 \in \Omega^1(M, \End E)$. The family $\nabla_t = \nabla_0 + t\alpha$ has curvature $F_t = F_0 + t\,d_{\nabla_0}\alpha + t^2\alpha \wedge \alpha$, so
$\frac{d}{dt}F_t = d_{\nabla_t}\alpha$. Hence
\[
\frac{d}{dt}P(F_t) = k\,\tilde{P}(d_{\nabla_t}\alpha, F_t, \ldots, F_t) = k\,d\,\tilde{P}(\alpha, F_t, \ldots, F_t),
\]
using the formula of (1) and $d_{\nabla_t}F_t = 0$. Integrating in $t$ from $0$ to $1$ gives $P(F_1) - P(F_0) = d\left(k\int_0^1\tilde{P}(\alpha, F_t, \ldots, F_t)\,dt\right)$, an exact form.

(3) The pullback connection has curvature $f^*F$, and $P$ commutes with pullback.
\end{proof}

\begin{definition}\label{characteristic-classes:definition-chern-forms}
The \emph{Chern forms} of $(E, \nabla)$ are defined by
\[
\det\left(I + \frac{i}{2\pi}F\right) = 1 + c_1(E, \nabla) + c_2(E, \nabla) + \cdots, \qquad c_k(E, \nabla) \in \Omega^{2k}(M),
\]
and the \emph{Chern classes} in de Rham cohomology are $c_k(E) = [c_k(E, \nabla)] \in H^{2k}_{\mathrm{dR}}(M)$, independent of $\nabla$ by Theorem~\ref{characteristic-classes:theorem-chern-weil}. The
\emph{total Chern class} is $c(E) = \sum_kc_k(E)$.
\end{definition}

Expanding the determinant with Example~\ref{characteristic-classes:example-invariant-polynomials}(2), the first Chern forms are
\[
c_1(E, \nabla) = \frac{i}{2\pi}\operatorname{tr}F, \qquad c_2(E, \nabla) = \frac{1}{8\pi^2}\bigl(\operatorname{tr}(F \wedge F) - \operatorname{tr}F \wedge \operatorname{tr}F\bigr), \qquad
c_r(E, \nabla) = \det\Bigl(\frac{i}{2\pi}F\Bigr),
\]
using $\sigma_2(A) = \frac12((\operatorname{tr}A)^2 - \operatorname{tr}A^2)$ with $A = \frac{i}{2\pi}F$ and $(\frac{i}{2\pi})^2 = -\frac{1}{4\pi^2}$. For a line bundle $F$ is a single $2$-form, and
$c_1 = \frac{i}{2\pi}F$ is the only Chern form. The factor $\frac{i}{2\pi}$ makes the classes real (for a Hermitian metric and a metric connection $F$ is skew-Hermitian, so
$\frac{i}{2\pi}F$ is Hermitian) and integral (Proposition~\ref{characteristic-classes:proposition-chern-properties}(3)).

\begin{example}\label{characteristic-classes:example-chern-forms-sphere}
\begin{enumerate}
\item \emph{The tautological bundle of $\CC\PP^1$.} With the projection connection, $F = 2i\,dx \wedge dy/(1 + |z|^2)^2$ in the chart $z = x + iy$
(Example~\ref{connections:example-curvature-computations}(3)), so $c_1(\mathcal{O}(-1), \nabla) = -\frac1\pi\,dx \wedge dy/(1 + |z|^2)^2$, and $\int_{\CC\PP^1}c_1 = -1$.
\item \emph{The tangent bundle of $S^2$.} The rotation by $+\pi/2$ in each tangent plane makes $TS^2$ a complex line bundle, with $i$ acting as this rotation. The projection
connection commutes with the rotation, and its curvature is $dA \otimes J$ with $J$ the rotation by $-\pi/2$ (Example~\ref{connections:example-curvature-computations}(2)), that is
$F = -i\,dA$ as a complex $1 \times 1$ matrix. So $c_1(TS^2, \nabla) = \frac{i}{2\pi}(-i)\,dA = \frac{1}{2\pi}dA$ and $\int_{S^2}c_1 = \frac{4\pi}{2\pi} = 2$. This is the degree of
$T\CC\PP^1$ found in Example~\ref{vector-bundles:example-degrees}(2), and the Euler characteristic of $S^2$: the Gauss--Bonnet theorem.
\item \emph{Flat and trivial bundles.} A flat connection has $F = 0$, so all Chern forms vanish. On a trivial bundle the connection $d + A$ has $F = dA + A \wedge A$ and
$c_1 = \frac{i}{2\pi}\operatorname{tr}(dA + A \wedge A) = \frac{i}{2\pi}d\operatorname{tr}A$, since $\operatorname{tr}(A \wedge A) = 0$ (the entries of $A \wedge A$ on the diagonal are
$\sum_kA^i_k \wedge A^k_i$, and the terms for $(i, k)$ and $(k, i)$ cancel); so $c_1$ is exact, as it must be.
\end{enumerate}
\end{example}

\begin{proposition}\label{characteristic-classes:proposition-chern-properties}
\begin{enumerate}
\item $c_0 = 1$ and $c_k(E) = 0$ for $k > r$; $c_k(f^*E) = f^*c_k(E)$.
\item $c(E \oplus E') = c(E)c(E')$ (the \emph{Whitney formula}).
\item For a line bundle $L$ with Hermitian metric and Chern connection, $c_1(L) = [\frac{i}{2\pi}F]$, which is the image in $H^2(M; \RR)$ of the integral class of
Theorem~\ref{vector-bundles:theorem-line-bundles}.
\item $c_k(\overline{E}) = (-1)^kc_k(E)$, and for the dual $c_k(E^\vee) = (-1)^kc_k(E)$.
\end{enumerate}
\end{proposition}

\begin{proof}
(1) The coefficients of $\det(I + tA)$ are the elementary symmetric functions of the eigenvalues, so they vanish above degree $r$; naturality is
Theorem~\ref{characteristic-classes:theorem-chern-weil}(3). (2) Choose connections $\nabla$, $\nabla'$; their direct sum has block diagonal curvature, and the determinant of a block diagonal matrix is the
product of the determinants. (4) A conjugate-linear isomorphism $E \cong \overline{E}$ carries a connection to the conjugate one with curvature $\overline{F}$, and
$\det(I + \frac{i}{2\pi}\overline{F}) = \overline{\det(I - \frac{i}{2\pi}F)}$; comparing degrees gives the sign, and $E^\vee \cong \overline{E}$ for a Hermitian metric
(Proposition~\ref{vector-bundles:proposition-metrics}).

(3) For a line bundle the curvature is a $2$-form $F$ and $c_1(L, \nabla) = \frac{i}{2\pi}F$. To identify the class, represent $L$ by a cocycle $(g_{ij})$ and choose a partition of unity; the \v{C}ech--de Rham
comparison (\cite[Chapter II]{bott-tu}) identifies the image of the integral class under $H^2(M; \ZZ) \to H^2_{\mathrm{dR}}(M)$ with the class of the $2$-form obtained from
$\frac{1}{2\pi i}d\log g_{ij}$, which is exactly $\frac{i}{2\pi}F$ for a connection built from the same trivialisations; see \cite[Chapter 1]{griffiths-harris}.
\end{proof}

\begin{example}\label{characteristic-classes:example-tautological}
For the tautological bundle $\mathcal{O}(-1)$ on $\CC\PP^1 = S^2$ with the metric induced from $\CC^2$, the Chern connection has curvature $F = -\frac{dz \wedge d\bar z}{(1 + |z|^2)^2}$ in the chart
$z$, and $\int_{\CC\PP^1}\frac{i}{2\pi}F = -1$; so $c_1(\mathcal{O}(-1))$ is $-1$ times the positive generator of $H^2(\CC\PP^1; \ZZ) \cong \ZZ$. Consequently $c_1(\mathcal{O}(1)) = h$ is the positive
generator, and $c(T\CC\PP^n) = (1 + h)^{n+1}$ by the Euler sequence (Example~\ref{vector-bundles:example-euler-sequence}) and the Whitney formula; for instance $c_1(T\CC\PP^n) = (n+1)h$.
\end{example}

\section{Chern classes: axioms and the splitting principle}\label{characteristic-classes:section-axioms}

\begin{theorem}\label{characteristic-classes:theorem-chern-axioms}
There is a unique way to assign to every complex vector bundle $E$ over a paracompact space classes $c_k(E) \in H^{2k}(M; \ZZ)$ such that: $c_0 = 1$; $c_k(f^*E) = f^*c_k(E)$; $c(E \oplus E') = c(E)c(E')$;
and $c_1(\mathcal{O}(-1))$ is the negative generator of $H^2(\CC\PP^1; \ZZ)$. Their images in de Rham cohomology are the Chern--Weil classes of
Definition~\ref{characteristic-classes:definition-chern-forms}.
\end{theorem}

This is \cite[Chapter 14]{milnor-stasheff}, where the classes are constructed from the Thom isomorphism or from the cohomology of projective bundles; the comparison with Chern--Weil theory is
\cite[Appendix C]{milnor-stasheff} or \cite[Chapter 1]{griffiths-harris}.

\begin{theorem}[Splitting principle]\label{characteristic-classes:theorem-splitting}
For every complex vector bundle $E \to M$ of rank $r$ there is a space $F(E)$ (the flag bundle) and a map $\pi \colon F(E) \to M$ such that $\pi^*$ is injective on cohomology and $\pi^*E$ is a direct sum of line
bundles $L_1 \oplus \cdots \oplus L_r$. Writing $x_i = c_1(L_i)$, the \emph{Chern roots}, one has $c(E) = \prod_i(1 + x_i)$, and every symmetric polynomial identity in the $x_i$ gives an identity between
characteristic classes.
\end{theorem}

This is \cite[Chapter 21]{bott-tu} or \cite[Chapter 14]{milnor-stasheff}: one takes the bundle of full flags in the fibres, whose cohomology is a free module over $H^\bullet(M)$ by the Leray--Hirsch
theorem, and the pullback of $E$ acquires a filtration with line bundle quotients, which splits after a further pullback or, over $\RR$, by choosing a metric.

\begin{definition}\label{characteristic-classes:definition-chern-character}
With Chern roots as above, the \emph{Chern character} and the \emph{Todd class} are
\[
\mathrm{ch}(E) = \sum_ie^{x_i} = r + c_1 + \tfrac12(c_1^2 - 2c_2) + \cdots, \qquad \mathrm{td}(E) = \prod_i\frac{x_i}{1 - e^{-x_i}} = 1 + \tfrac12c_1 + \tfrac{1}{12}(c_1^2 + c_2) + \cdots,
\]
elements of $H^\bullet(M; \QQ)$, well defined because they are symmetric power series in the roots, hence polynomials in the Chern classes.
\end{definition}

\begin{proposition}\label{characteristic-classes:proposition-chern-character}
$\mathrm{ch}(E \oplus E') = \mathrm{ch}(E) + \mathrm{ch}(E')$ and $\mathrm{ch}(E \otimes E') = \mathrm{ch}(E)\mathrm{ch}(E')$; the Todd class is multiplicative in short exact sequences. Hence $\mathrm{ch}$ induces a
ring homomorphism from the K-theory of $M$ to $H^{\mathrm{even}}(M; \QQ)$.
\end{proposition}

\begin{proof}
By the splitting principle we may assume $E = \bigoplus L_i$ and $E' = \bigoplus L'_j$; then the Chern roots of $E \oplus E'$ are the $x_i$ and $y_j$, and those of $E \otimes E'$ are the $x_i + y_j$, because
$c_1(L \otimes L') = c_1(L) + c_1(L')$ (Theorem~\ref{vector-bundles:theorem-line-bundles}). The statements follow from $e^{x+y} = e^xe^y$ and from multiplicativity of products over the roots. Short exact
sequences of bundles split topologically (Proposition~\ref{vector-bundles:proposition-metrics}), which gives the last statements.
\end{proof}

\begin{example}\label{characteristic-classes:example-chern-character-cp1}
On $\CC\PP^1$ let $h \in H^2(\CC\PP^1; \ZZ)$ be the positive generator, so $\int h = 1$ and $h^2 = 0$. Then $c_1(\mathcal{O}(k)) = kh$, $\mathrm{ch}(\mathcal{O}(k)) = e^{kh} = 1 + kh$, and
$T\CC\PP^1 \cong \mathcal{O}(2)$ has $\mathrm{td}(T\CC\PP^1) = 1 + \frac12c_1 = 1 + h$. The Riemann--Roch formula (Remark~\ref{characteristic-classes:remark-signature}) then predicts
$\chi(\CC\PP^1, \mathcal{O}(k)) = \int(1 + kh)(1 + h) = k + 1$, which is the number of homogeneous polynomials of degree $k$ in two variables when $k \geq 0$, the dimension of
the space of holomorphic sections of $\mathcal{O}(k)$ (Chapter~\ref{holomorphic-bundles}).
\end{example}

\section{The Thom isomorphism in de Rham cohomology}\label{characteristic-classes:section-thom}

In this section $\pi \colon E \to M$ is a smooth real vector bundle of rank $n$ over a manifold $M$ of dimension $m$. It is \emph{oriented} if it has an atlas of local trivialisations $E|_U \cong U \times \RR^n$ whose transition
functions take values in $\GL_n^+(\RR)$, the matrices of positive determinant; such trivialisations are called oriented.

\begin{definition}\label{characteristic-classes:definition-vertical-support}
A form $\omega$ on $E$ has \emph{compact vertical support} if $\operatorname{supp}\omega \cap \pi^{-1}(K)$ is compact for every compact $K \subseteq M$. These forms form a subcomplex $\Omega^\bullet_{\mathrm{cv}}(E)$ of $\Omega^\bullet(E)$,
whose cohomology is written $H^\bullet_{\mathrm{cv}}(E)$. Restriction to $E|_U$ for $U \subseteq M$ open preserves compact vertical support.
\end{definition}

Let $t = (t_1, \ldots, t_n)$ be the fibre coordinates of an oriented trivialisation $E|_U \cong U \times \RR^n$. Every $\omega \in \Omega^k_{\mathrm{cv}}(E|_U)$ is uniquely a sum $\omega = \omega' + \omega''$, where $\omega''$ is the part containing
$dt_1 \wedge \cdots \wedge dt_n$ and $\omega'$ contains fewer than $n$ of the $dt_j$. Locally on $U$, in coordinates $x$ of $M$, $\omega'' = \sum_If_I(x, t)\,dx^I \wedge dt_1 \wedge \cdots \wedge dt_n$.

\begin{definition}\label{characteristic-classes:definition-fibre-integration}
For an oriented bundle, \emph{integration along the fibres} $\pi_* \colon \Omega^k_{\mathrm{cv}}(E) \to \Omega^{k-n}(M)$ is defined in an oriented trivialisation by $\pi_*\omega' = 0$ and
\[
\pi_*\Big(\sum_If_I(x, t)\,dx^I \wedge dt_1 \wedge \cdots \wedge dt_n\Big) = \sum_I\Big(\int_{\RR^n}f_I(x, t)\,dt\Big)dx^I.
\]
\end{definition}

\begin{proposition}\label{characteristic-classes:proposition-fibre-integration}
Integration along the fibres is well defined, commutes with restriction to open subsets of $M$ and with $d$, $d\pi_* = \pi_*d$, and satisfies the \emph{projection formula} $\pi_*(\pi^*\alpha \wedge \omega) = \alpha \wedge \pi_*\omega$ for
$\alpha \in \Omega^\bullet(M)$ and $\omega \in \Omega^\bullet_{\mathrm{cv}}(E)$. Reversing the orientation of $E$ changes the sign of $\pi_*$.
\end{proposition}

\begin{proof}
The integrals converge, since $f_I(x, -)$ has compact support, locally uniformly in $x$. A change of oriented trivialisation $t' = g(x)t$ with $\det g > 0$ gives $dt'_1 \wedge \cdots \wedge dt'_n = \det g(x)\,dt_1 \wedge \cdots \wedge dt_n$
plus terms containing some $dx^i$ and fewer $dt_j$; so the parts $\omega'$ and $\omega''$ are the same in both trivialisations, and the integrals agree by the change of variables formula
(Theorem~\ref{measure-theory:theorem-change-of-variables}), with Jacobian $\det g(x) = |\det g(x)|$. A change of coordinates in $M$ does not affect the fibre integrals. So $\pi_*$ is well defined, and it clearly commutes with
restrictions. For the projection formula, $\pi^*\alpha \wedge \omega'$ has fewer than $n$ factors $dt_j$, and $\pi^*(dx^K) \wedge f\,dx^I \wedge dt_1 \wedge \cdots \wedge dt_n = f\,dx^K \wedge dx^I \wedge dt_1 \wedge \cdots \wedge dt_n$; by linearity
in $\alpha$ over $C^\infty(M)$ this gives the formula. For $d$: if $\omega = f\,dx^I \wedge dt_J$ with $|J| < n$, the part of $d\omega$ with $n$ factors $dt$ is a sum of terms $\pm\partial_{t_j}f\,dx^I \wedge dt_1 \wedge \cdots \wedge dt_n$, whose fibre
integrals vanish because $\int_\RR\partial_{t_j}f\,dt_j = 0$ for compactly supported $f$ (Theorem~\ref{real-analysis:theorem-ftc}); so $\pi_*d\omega = 0 = d\pi_*\omega$. If $\omega = f\,dx^I \wedge dt_1 \wedge \cdots \wedge dt_n$, then
$d\omega = \sum_i\partial_{x_i}f\,dx^i \wedge dx^I \wedge dt_1 \wedge \cdots \wedge dt_n$, and $\pi_*d\omega = \sum_i(\int\partial_{x_i}f\,dt)\,dx^i \wedge dx^I = d\pi_*\omega$ by differentiation under the integral sign
(Proposition~\ref{real-analysis:proposition-parameter-integrals}). Reversing the orientation replaces $dt_1 \wedge \cdots \wedge dt_n$ by $-dt_1 \wedge \cdots \wedge dt_n$ in an oriented trivialisation.
\end{proof}

\begin{lemma}\label{characteristic-classes:lemma-thom-trivial}
For every manifold $N$ and $n \geq 0$, integration along the fibres of the trivial bundle $N \times \RR^n \to N$ induces isomorphisms $H^k_{\mathrm{cv}}(N \times \RR^n) \to H^{k-n}(N)$.
\end{lemma}

\begin{proof}
We treat $n = 1$ in a form that allows induction. Let $N' = N \times \RR^{n-1}$ and consider the projection $p \colon N' \times \RR \to N'$ with fibre coordinate $t$, and the subcomplexes $C^\bullet \subseteq \Omega^\bullet(N' \times \RR)$ and
$C'^\bullet \subseteq \Omega^\bullet(N')$ of forms whose supports meet the preimage of every compact subset of $N$ in a compact set. Every $\omega \in C^k$ is uniquely $\omega = \alpha_t + \beta_t \wedge dt$, with $\alpha_t$, $\beta_t$ families of
forms on $N'$ depending smoothly on $t$ and without $dt$, and $p_*\omega = \int_\RR\beta_t\,dt$, which lies in $C'^{k-1}$. Choose $e \in C^\infty_c(\RR)$ with $\int e = 1$, let $E(t) = \int_{-\infty}^te$, and define
\[
e_*(\beta) = p^*\beta \wedge e(t)\,dt, \qquad K\omega = \int_{-\infty}^t\beta_u\,du - E(t)\int_{-\infty}^\infty\beta_u\,du.
\]
Then $e_* \colon C'^{k-1} \to C^k$ commutes with $d$ (as $d(e\,dt) = 0$), $p_*e_* = \mathrm{id}$, and $K\omega$ lies in $C^{k-1}$: it vanishes for $t$ below the support of $\omega$, and also above it, where $\int_{-\infty}^t\beta_u\,du = \int_{-\infty}^\infty\beta_u\,du$
and $E(t) = 1$ (using that the $t$-supports are uniformly bounded over compact subsets of $N$). We claim that on $C^k$,
\[
dK\omega - Kd\omega = (-1)^{k-1}(\omega - e_*p_*\omega).
\]
If $\omega = \alpha_t$ has no $dt$, then $K\omega = 0$, $e_*p_*\omega = 0$, and the $dt$-part of $d\omega$ is $(-1)^k\partial_t\alpha_t \wedge dt$, so $Kd\omega = (-1)^k\int_{-\infty}^t\partial_u\alpha_u\,du = (-1)^k\alpha_t$ (the second term vanishes as
$\int_\RR\partial_u\alpha_u\,du = 0$), which gives the claim. If $\omega = \beta_t \wedge dt$, write $d_{N'}$ for the exterior derivative along $N'$ and $g_t = K\omega$. Then $d\omega = d_{N'}\beta_t \wedge dt$, so
$Kd\omega = d_{N'}g_t$ by differentiation under the integral sign, while $dK\omega = d_{N'}g_t + (-1)^{k-1}\partial_tg_t \wedge dt$ with $\partial_tg_t = \beta_t - e(t)\int_\RR\beta_u\,du$. Hence
$dK\omega - Kd\omega = (-1)^{k-1}(\beta_t \wedge dt - e(t)(\int_\RR\beta_u\,du) \wedge dt) = (-1)^{k-1}(\omega - e_*p_*\omega)$, proving the claim. So $p_*$ and $e_*$ are inverse isomorphisms on cohomology.

For $n = 1$ this is the lemma. For general $n$, integration over $\RR^n$ is the composite of integration over the last coordinate, which maps the complex of forms on $N \times \RR^n$ with compact vertical support over $N$ to
the corresponding complex on $N \times \RR^{n-1}$, and integration over $\RR^{n-1}$; by Fubini's theorem (in the form of Proposition~\ref{real-analysis:proposition-parameter-integrals}) and the orientation convention
$dt_1 \wedge \cdots \wedge dt_n = (dt_1 \wedge \cdots \wedge dt_{n-1}) \wedge dt_n$, and induction on $n$, the lemma follows.
\end{proof}

\begin{theorem}[Thom isomorphism]\label{characteristic-classes:theorem-thom-de-rham}
Let $\pi \colon E \to M$ be an oriented vector bundle of rank $n$ over a manifold with a finite good cover (for instance a compact manifold). Then $\pi_* \colon H^k_{\mathrm{cv}}(E) \to H^{k-n}(M)$ is an isomorphism for every $k$. There is a
unique class $\Phi \in H^n_{\mathrm{cv}}(E)$ with $\pi_*\Phi = 1$, the \emph{Thom class}, and the inverse isomorphism is $\alpha \mapsto \pi^*\alpha \wedge \Phi$.
\end{theorem}

\begin{proof}
For open $U, V \subseteq M$ and a partition of unity $\{\rho_U, \rho_V\}$ subordinate to $\{U, V\}$ on $U \cup V$, the sequence
$0 \to \Omega_{\mathrm{cv}}(E|_{U \cup V}) \to \Omega_{\mathrm{cv}}(E|_U) \oplus \Omega_{\mathrm{cv}}(E|_V) \to \Omega_{\mathrm{cv}}(E|_{U \cap V}) \to 0$ is exact, exactly as in Theorem~\ref{de-rham:theorem-mayer-vietoris}, using the functions
$\pi^*\rho_U$, $\pi^*\rho_V$: multiplying a form on $E|_{U \cap V}$ with compact vertical support by $\pi^*\rho_V$ and extending by zero gives a form on $E|_U$ with compact vertical support, because a compact $K \subseteq U$ meets
$\operatorname{supp}\rho_V$ in a compact subset of $U \cap V$. By Proposition~\ref{characteristic-classes:proposition-fibre-integration}, $\pi_*$ maps this sequence to the Mayer--Vietoris sequence of $U$, $V$ in de Rham
cohomology, commuting with restrictions, and it commutes with the connecting maps, which are given on both sides by $\omega \mapsto d(\pi^*\rho_V\,\omega)$ and $\beta \mapsto d(\rho_V\beta)$, since
$\pi_*d(\pi^*\rho_V\,\omega) = d(\rho_V\pi_*\omega)$ by the projection formula. So the long exact sequences form a commutative ladder.

Choose a finite good cover $\{U_1, \ldots, U_N\}$ of $M$ subordinate to a cover by trivialising open sets (possible by \cite[Chapter I, Section 5]{bott-tu}): all nonempty finite intersections of the $U_i$ are diffeomorphic to $\RR^m$ and
$E$ is trivial over each. We show by induction on $k$ that $\pi_*$ is an isomorphism over any union of $k$ open sets forming such a good cover. For $k = 1$ this is Lemma~\ref{characteristic-classes:lemma-thom-trivial}. For the
induction step write the union as $W \cup U$ with $W$ a union of $k - 1$ of the sets; then $W \cap U$ is the union of the $k - 1$ sets $U_i \cap U$, which again form such a cover. By induction $\pi_*$ is an isomorphism over $W$, $U$ and
$W \cap U$, and by the five lemma (Lemma~\ref{homological-algebra:lemma-five}) over $W \cup U$.

The class $\Phi = \pi_*^{-1}(1)$ is unique. By the projection formula $\pi_*(\pi^*\alpha \wedge \Phi) = \alpha \wedge \pi_*\Phi = \alpha$, so $\alpha \mapsto \pi^*\alpha \wedge \Phi$ is the inverse of $\pi_*$.
\end{proof}

\begin{definition}\label{characteristic-classes:definition-euler-de-rham}
For $E$ and $M$ as in Theorem~\ref{characteristic-classes:theorem-thom-de-rham}, the \emph{Euler class} of $E$ in de Rham cohomology is $e(E) = s^*\Phi \in H^n(M)$, where $s \colon M \to E$ is the zero section and $\Phi$ is viewed as a class in
$H^n(E)$.
\end{definition}

\begin{proposition}\label{characteristic-classes:proposition-euler-de-rham}
Let $M$ be compact and $E$, $F$ oriented bundles over $M$ of ranks $n$, $n'$.
\begin{enumerate}
\item $e(f^*E) = f^*e(E)$ for a smooth map $f \colon N \to M$ of compact manifolds, and reversing the orientation of $E$ changes the sign of $e(E)$.
\item If $n$ is odd, $e(E) = 0$.
\item If $E$ has a nowhere vanishing section, $e(E) = 0$.
\item $e(E \oplus F) = e(E) \wedge e(F)$, with $E \oplus F$ oriented by an oriented frame of $E$ followed by one of $F$.
\end{enumerate}
\end{proposition}

\begin{proof}
We use that a closed form $\phi \in \Omega^n_{\mathrm{cv}}(E)$ represents $\Phi$ if and only if $\pi_*\phi = 1$: $\pi_*\phi$ is a closed $0$-form, that is, locally constant, and it represents $\pi_*[\phi]$. (1) The bundle map
$\tilde f \colon f^*E \to E$ is an orientation-preserving isomorphism on fibres, so $\tilde f^*\phi$ has compact vertical support and, computing in pulled back trivialisations, fibre integral $f^*\pi_*\phi = 1$; hence it represents
the Thom class of $f^*E$, and $s_N^*\tilde f^*\phi = f^*s^*\phi$. The orientation statement follows from the last statement of Proposition~\ref{characteristic-classes:proposition-fibre-integration}.
(2) Let $a(v) = -v$. In oriented trivialisations the change of variables $t \mapsto -t$ has Jacobian $(-1)^n = -1$, so $\pi_*a^*\phi = -\pi_*\phi = -1$, and $a^*\phi$ represents $-\Phi$. As $a \circ s = s$,
$e(E) = [s^*\phi] = [s^*a^*\phi] = -e(E)$.
(3) Choose a metric on $E$ and let $\sigma$ be a section with $|\sigma| \geq c > 0$ (by compactness). For $\lambda > 0$ let $m_\lambda(v) = \lambda v$; then $m_\lambda^*\phi$ has fibre integral $1$ (Jacobian $\lambda^n > 0$), so it represents $\Phi$, and its support is
$\lambda^{-1}\operatorname{supp}\phi$. As $\operatorname{supp}\phi$ lies in some $\{|v| \leq R\}$, for $\lambda > R/c$ the form $m_\lambda^*\phi$ vanishes along $\sigma(M)$, so $\sigma^*m_\lambda^*\phi = 0$. The sections $s$ and $\sigma$ are homotopic
through $t\sigma$, so $e(E) = s^*[m_\lambda^*\phi] = \sigma^*[m_\lambda^*\phi] = 0$ by homotopy invariance (Theorem~\ref{de-rham:theorem-homotopy-invariance}).
(4) Let $p_1 \colon E \oplus F \to E$ and $p_2 \colon E \oplus F \to F$ be the projections and $\phi$, $\psi$ Thom forms. Then $p_1^*\phi \wedge p_2^*\psi$ has compact vertical support, and in product trivialisations with fibre coordinates
$(t, u)$ its component containing $dt_1 \wedge \cdots \wedge dt_n \wedge du_1 \wedge \cdots \wedge du_{n'}$ is the product of the components of $\phi$ and $\psi$ of top vertical degree, so by Fubini its fibre integral is
$\pi_*\phi \cdot \pi_*\psi = 1$. Hence it represents the Thom class of $E \oplus F$, and pulling back by the zero section, which is $(s_E, s_F)$, gives $e(E \oplus F) = e(E) \wedge e(F)$.
\end{proof}

\begin{remark}\label{characteristic-classes:remark-thom-poincare-dual}
The Thom class links Poincar\'e duality and characteristic classes. If $S \subseteq M$ is a compact oriented submanifold of a compact oriented manifold, with normal bundle $\nu$ identified with a tubular neighbourhood of $S$, the
Thom form of $\nu$, extended by zero, represents the Poincar\'e dual of $S$ in $M$; restricted to $S$ it gives $e(\nu)$, the \emph{self-intersection} of $S$. Applied to the diagonal of $M \times M$, whose normal bundle is $TM$, this
yields $\int_Me(TM) = \chi(M)$, part (3) of Theorem~\ref{characteristic-classes:theorem-euler} \cite{bott-tu}.
\end{remark}

\section{Euler, Pontryagin and Stiefel--Whitney classes}\label{characteristic-classes:section-euler}

\begin{theorem}\label{characteristic-classes:theorem-euler}
Let $E \to M$ be an oriented real vector bundle of rank $n$ over a compact manifold.
\begin{enumerate}
\item (Thom isomorphism) $H^{k}(M; \ZZ) \cong H^{k+n}(E, E \setminus 0; \ZZ)$, given by cup product with a \emph{Thom class} $u \in H^n(E, E \setminus 0; \ZZ)$ restricting to a generator on each fibre.
\item The \emph{Euler class} $e(E) \in H^n(M; \ZZ)$ is the restriction of $u$ to the zero section. It is natural, satisfies $e(E \oplus E') = e(E)e(E')$, vanishes if $E$ has a nowhere vanishing section, and for
a complex bundle of rank $r$ one has $e(E) = c_r(E)$.
\item (Gauss--Bonnet--Chern and Poincar\'e--Hopf) For a compact oriented $n$-manifold, $\int_Me(TM) = \chi(M)$, and the sum of the indices of the zeros of a vector field with isolated zeros equals $\chi(M)$.
\end{enumerate}
\end{theorem}

These are \cite[Chapters 9--11]{milnor-stasheff} and \cite[Chapter 6]{bott-tu}. With real coefficients, (1) is Theorem~\ref{characteristic-classes:theorem-thom-de-rham} (whose proof works over
any manifold with a finite good cover), and the de Rham image of the Euler class is the class of Definition~\ref{characteristic-classes:definition-euler-de-rham}, whose
properties were proved in Proposition~\ref{characteristic-classes:proposition-euler-de-rham}. In the Chern--Weil picture, the Pfaffian of the curvature of a metric connection, divided by
$(2\pi)^{n/2}$, is a closed form whose integral over $M$ is computed in Section~\ref{characteristic-classes:section-gauss-bonnet-chern} by Chern's intrinsic argument: it is the sum of the
indices of any vector field with isolated zeros, hence $\chi(M)$ by the Poincar\'e--Hopf theorem in (3).

\begin{definition}\label{characteristic-classes:definition-pontryagin}
For a real vector bundle $E$, the \emph{Pontryagin classes} are $p_k(E) = (-1)^kc_{2k}(E \otimes \CC) \in H^{4k}(M; \ZZ)$, and the \emph{Stiefel--Whitney classes} $w_k(E) \in H^k(M; \ZZ/2\ZZ)$ are the mod $2$
characteristic classes determined by the same axioms with $\ZZ/2$ coefficients; $w_1(E) = 0$ if and only if $E$ is orientable.
\end{definition}

\begin{remark}\label{characteristic-classes:remark-signature}
Hirzebruch's signature theorem expresses the signature of a compact oriented $4m$-manifold (Definition~\ref{poincare-duality:definition-intersection-form}) as $\int_ML(p_1, \ldots, p_m)$ for a universal
polynomial in the Pontryagin classes, and the Riemann--Roch--Hirzebruch theorem computes the holomorphic Euler characteristic of a holomorphic bundle on a compact complex manifold as
$\int_X\mathrm{ch}(E)\,\mathrm{td}(TX)$ \cite{hirzebruch-topological}. Both are special cases of the Atiyah--Singer index theorem (Chapter~\ref{index-theorems}), and both are used in Chapter~\ref{kodaira} and
Chapter~\ref{intersection-theory}.
\end{remark}

\begin{remark}\label{characteristic-classes:remark-classifying-spaces}
Isomorphism classes of rank $r$ complex bundles over a paracompact space $M$ correspond to homotopy classes of maps $M \to \mathrm{BU}(r)$, where $\mathrm{BU}(r)$ is the infinite Grassmannian
$\mathrm{Gr}_r(\CC^\infty)$, and $H^\bullet(\mathrm{BU}(r); \ZZ) = \ZZ[c_1, \ldots, c_r]$ is a polynomial ring on the universal Chern classes; every characteristic class is a pullback of a universal one
\cite[Chapter 14]{milnor-stasheff}. For line bundles, $\mathrm{BU}(1) = \CC\PP^\infty = K(\ZZ, 2)$, which recovers Theorem~\ref{vector-bundles:theorem-line-bundles} and
Theorem~\ref{cohomology-operations:theorem-representability}.
\end{remark}

\begin{example}\label{characteristic-classes:example-algebraic-chern}
On a smooth projective complex variety, the Chern classes of an algebraic vector bundle are algebraic cycle classes, hence Hodge classes; for a line bundle, $c_1(L)$ is the class of a divisor
(Chapter~\ref{divisors}), and the Lefschetz theorem on $(1,1)$-classes (Chapter~\ref{lefschetz-11}) says conversely that every integral $(1,1)$-class is of this form. This is the case $p = 1$ of the Hodge
conjecture, and it is the only case known in general.
\end{example}

\section{The Euler form and the Gauss--Bonnet--Chern theorem}\label{characteristic-classes:section-gauss-bonnet-chern}

Chern--Weil theory for complex bundles uses polynomials that are invariant under all of $\GL_r(\CC)$. A real oriented vector bundle with a metric has more structure: in oriented
orthonormal frames the curvature of a metric connection is a skew-symmetric matrix of $2$-forms, and the changes of frame take values in $\mathrm{SO}(n)$. Skew-symmetric matrices have
an invariant that is not a polynomial in the coefficients of the characteristic polynomial, the Pfaffian (Definition~\ref{multilinear-algebra:definition-pfaffian}), a square root of the
determinant. Applied to the curvature it gives a closed form, the \emph{Euler form}. This section proves, following Chern \cite{chern-1944}, that when the rank of the bundle equals the
dimension of the base, the integral of the Euler form counts the zeros of any section, with signs:
\[
\int_Me(\nabla) = \sum_{s(z) = 0}\operatorname{ind}_z(s) .
\]
For the tangent bundle the right side is the Euler characteristic by the Poincar\'e--Hopf theorem, and the formula becomes the \emph{Gauss--Bonnet--Chern theorem}
$\int_Me(TM, \nabla) = \chi(M)$, the higher-dimensional Gauss--Bonnet theorem. Chern's proof is \emph{intrinsic}: earlier proofs, by Allendoerfer, Fenchel and Weil, used isometric embeddings of $M$, or of
pieces of it, into Euclidean spaces, while Chern only computes on the unit sphere bundle of $E$, where the pullback of the Euler form becomes exact, with an explicit primitive. The surface case, in which the whole argument fits in a few
lines, is worked out in Example~\ref{characteristic-classes:example-chern-surfaces}, and it may help to read it first.

\emph{On a first reading} one can read the definitions and the statements, the example of surfaces, and Corollary~\ref{characteristic-classes:corollary-gauss-bonnet-chern}; the
transgression computation (Lemma~\ref{characteristic-classes:lemma-transgression} and Proposition~\ref{characteristic-classes:proposition-transgression}) is not used elsewhere.

Throughout, $E \to M$ is an oriented real vector bundle of rank $n$ with a metric $\langle\ ,\ \rangle$ on its fibres, and $\nabla$ is a metric connection
(Proposition~\ref{connections:proposition-metric-connections}). Near every point there are \emph{oriented orthonormal frames} $e = (e_1, \ldots, e_n)$, obtained from any oriented frame by
the Gram--Schmidt process. Since such a frame is orthonormal we write its connection form and curvature with lower indices, $\theta_{ij} = \theta^i_j$ and $F_{ij} = F^i_j$, so that
\[
\nabla e_j = \sum_i\theta_{ij} \otimes e_i, \qquad F_{ij} = d\theta_{ij} + \sum_k\theta_{ik} \wedge \theta_{kj}, \qquad dF = F \wedge \theta - \theta \wedge F
\]
(Propositions~\ref{connections:proposition-connections-exist} and~\ref{connections:proposition-curvature}). Both matrices are skew-symmetric, $\theta_{ji} = -\theta_{ij}$ and
$F_{ji} = -F_{ij}$, by Proposition~\ref{connections:proposition-metric-connections}. Another oriented orthonormal frame is $e^{\prime} = eg$ with $g \colon U \to \mathrm{SO}(n)$, and then
$\theta^{\prime} = g^{-1}\theta g + g^{-1}dg$ and $F^{\prime} = g^{-1}Fg = g^TFg$.

\begin{lemma}\label{characteristic-classes:lemma-normal-frames}
Let $x \in M$ and let $(b_1, \ldots, b_n)$ be an oriented orthonormal basis of $E_x$. There is an oriented orthonormal frame $e$ on a neighbourhood of $x$ with $e_i(x) = b_i$ and whose
connection form vanishes at $x$, $\theta(x) = 0$.
\end{lemma}

\begin{proof}
Let $f$ be an oriented orthonormal frame on a coordinate neighbourhood of $x$, with coordinates $x^1, \ldots, x^d$ vanishing at $x$, and connection form $\theta_f = \sum_k\theta_k\,dx^k$
with skew-symmetric matrices of functions $\theta_k$. Let $A_k = \theta_k(x)$ and let $h \in \mathrm{SO}(n)$ be the matrix with $f(x)h = (b_1, \ldots, b_n)$. Put
$g = \exp\bigl(-\sum_kx^kA_k\bigr)h$. The exponential of a skew-symmetric matrix $B$ is orthogonal, as $\exp(B)^T = \exp(-B) = \exp(B)^{-1}$, with determinant $e^{\operatorname{tr}B} = 1$; so
$g$ takes values in $\mathrm{SO}(n)$ and $e = fg$ is an oriented orthonormal frame with $e(x) = f(x)h = b$. At $x$ we have $g(x) = h$ and $dg(x) = -\sum_kA_k\,dx^k\,h = -\theta_f(x)h$,
so $\theta_e(x) = h^{-1}\theta_f(x)h + h^{-1}dg(x) = 0$.
\end{proof}

Such a frame is the analogue of normal coordinates: at the single point $x$ the connection looks like the trivial one, while the curvature, which is a tensor, cannot be removed. It
simplifies computations of quantities that do not depend on the frame: the second Bianchi identity says that $dF(x) = 0$ in such a frame.

Let $n = 2m$ be even. Since $2$-forms commute, the Pfaffian can be applied to the skew-symmetric matrix of $2$-forms $F$, using the formula of
Proposition~\ref{multilinear-algebra:proposition-pfaffian}(1):
\[
\mathrm{Pf}(F) = \frac{1}{2^mm!}\sum_{\sigma \in S_n}\operatorname{sgn}(\sigma)\,F_{\sigma(1)\sigma(2)} \wedge F_{\sigma(3)\sigma(4)} \wedge \cdots \wedge F_{\sigma(n-1)\sigma(n)} .
\]

\begin{definition}\label{characteristic-classes:definition-euler-form}
For $n = 2m$, the \emph{Euler form} of $(E, \nabla)$ is the $n$-form
\[
e(\nabla) = \mathrm{Pf}\Bigl(\frac{F}{2\pi}\Bigr) = \frac{1}{(2\pi)^m}\mathrm{Pf}(F) \in \Omega^n(M),
\]
computed in any oriented orthonormal frame. For $n = 2$, $e(\nabla) = \frac{1}{2\pi}F_{12}$; for $n = 4$,
$e(\nabla) = \frac{1}{4\pi^2}(F_{12} \wedge F_{34} - F_{13} \wedge F_{24} + F_{14} \wedge F_{23})$.
\end{definition}

\begin{proposition}\label{characteristic-classes:proposition-euler-form}
Let $n = 2m$.
\begin{enumerate}
\item $e(\nabla)$ does not depend on the oriented orthonormal frame, so it is a global $n$-form, and reversing the orientation of $E$ changes its sign.
\item $e(\nabla)$ is closed.
\item Its class in $H^n_{\mathrm{dR}}(M)$ does not depend on the metric connection $\nabla$.
\item It is natural: $e(f^*\nabla) = f^*e(\nabla)$ for a smooth map $f \colon N \to M$.
\item If $E$ is the underlying oriented real bundle of a complex line bundle $L$ with a Hermitian metric $h$, with the metric $\operatorname{Re}h$, and $\nabla$ is a metric connection on $L$,
then $e(\nabla) = c_1(L, \nabla)$.
\end{enumerate}
\end{proposition}

\begin{proof}
(1) For $e^{\prime} = eg$ with $g$ in $\mathrm{SO}(n)$, $\mathrm{Pf}(F^{\prime}) = \mathrm{Pf}(g^TFg) = \det(g)\mathrm{Pf}(F) = \mathrm{Pf}(F)$ by
Proposition~\ref{multilinear-algebra:proposition-pfaffian}(2), which applies to matrices with entries in the commutative algebra of even forms at a point. Reversing the orientation
replaces oriented frames by frames $eg$ with $\det g = -1$.

(2) $\mathrm{Pf}(F)$ is a polynomial with constant coefficients in the $2$-forms $F_{ij}$, which commute, so by the Leibniz rule $d\,\mathrm{Pf}(F)$ is a sum of terms each containing one
factor $dF_{ij}$. At a point $x$ compute in a frame as in Lemma~\ref{characteristic-classes:lemma-normal-frames}: there $dF(x) = (F \wedge \theta - \theta \wedge F)(x) = 0$, so
$d\,\mathrm{Pf}(F)$ vanishes at $x$. As $x$ is arbitrary and $\mathrm{Pf}(F)$ does not depend on the frame, $d\,\mathrm{Pf}(F) = 0$.

(3) Let $\widetilde{\mathrm{Pf}}$ be the polarisation of $\mathrm{Pf}$, the symmetric $m$-linear form on skew-symmetric matrices with $\widetilde{\mathrm{Pf}}(A, \ldots, A) = \mathrm{Pf}(A)$; it is
given by a universal formula in terms of $\mathrm{Pf}$ (for instance $m!\,\widetilde{\mathrm{Pf}}(A_1, \ldots, A_m)$ is the coefficient of $t_1t_2 \cdots t_m$ in $\mathrm{Pf}(t_1A_1 + \cdots + t_mA_m)$), so it is
invariant under $A_i \mapsto g^TA_ig$ for $g \in \mathrm{SO}(n)$. Differentiating this invariance along $g = \exp(tB)$ with $B$ skew-symmetric, for which
$\frac{d}{dt}\big|_{t=0}\exp(tB)^TA\exp(tB) = AB - BA$, gives $\sum_i\widetilde{\mathrm{Pf}}(A_1, \ldots, [B, A_i], \ldots, A_m) = 0$ for skew-symmetric $B$ and $A_i$. The proof of
Lemma~\ref{characteristic-classes:lemma-invariant-derivative} uses only this infinitesimal invariance, for $B$ the components of the connection form, and these are skew-symmetric for a
metric connection in an orthonormal frame. So that lemma holds for $\widetilde{\mathrm{Pf}}$ and skew-symmetric matrices of forms, and the proof of
Theorem~\ref{characteristic-classes:theorem-chern-weil}(2) applies word for word: for metric connections $\nabla_0$, $\nabla_1$ the connections $\nabla_t = \nabla_0 + t(\nabla_1 - \nabla_0)$
are metric, the matrix of $\alpha = \nabla_1 - \nabla_0$ in an orthonormal frame is skew-symmetric, and
$\mathrm{Pf}(F_1) - \mathrm{Pf}(F_0) = d\bigl(m\int_0^1\widetilde{\mathrm{Pf}}(\alpha, F_t, \ldots, F_t)\,dt\bigr)$, where the form in brackets does not depend on the oriented orthonormal frame by the
invariance of $\widetilde{\mathrm{Pf}}$.

(4) The pullback of an oriented orthonormal frame is one for $f^*E$, with connection form $f^*\theta$ and curvature $f^*F$.

(5) Let $s$ be a local frame of $L$ with $h(s, s) = 1$. Then $(e_1, e_2) = (s, is)$ is an oriented orthonormal frame of $E$ for $\operatorname{Re}h$. The connection form of $\nabla$ on $L$ in
the frame $s$ is $i\omega$ for a real $1$-form $\omega$ (it is skew-Hermitian), so $\nabla e_1 = i\omega \otimes s = \omega \otimes e_2$ and $\nabla e_2 = i\nabla s = -\omega \otimes e_1$. Hence
$\theta_{21} = \omega$, $\theta_{12} = -\omega$, and $F_{12} = d\theta_{12} = -d\omega$ (for $n = 2$ the term $\theta \wedge \theta$ has zero off-diagonal entries). The curvature of $L$ is
$d(i\omega) = i\,d\omega$, so $c_1(L, \nabla) = \frac{i}{2\pi}\,i\,d\omega = -\frac{1}{2\pi}d\omega = \frac{1}{2\pi}F_{12} = e(\nabla)$.
\end{proof}

By (5) and Example~\ref{characteristic-classes:example-chern-forms-sphere}(2), for the round sphere $S^2$ the Euler form of the projection connection on $TS^2$ is $\frac{1}{2\pi}dA$, with
integral $2 = \chi(S^2)$. The rest of the section shows that this is no accident.

Next we define the index of an isolated zero of a section, which generalises the winding number of a vector field in the plane around a zero.

Write $\sigma$ for the $(n-1)$-form
\[
\sigma = \sum_{i=1}^n(-1)^{i-1}y^i\,dy^1 \wedge \cdots \wedge \widehat{dy^i} \wedge \cdots \wedge dy^n
\]
on $\RR^n$, and also for its restriction to the unit sphere $S^{n-1}$, oriented as the boundary of the unit ball. At $y$ it is the interior product of $dy^1 \wedge \cdots \wedge dy^n$ with
the position vector $y$: $\sigma_y(v_2, \ldots, v_n) = \det(y, v_2, \ldots, v_n)$. Hence $\sigma$ restricts to the area form of $S^{n-1}$, and for a linear map $h \in \mathrm{O}(n)$, which maps
$S^{n-1}$ to itself, $h^*\sigma = \det(h)\,\sigma$, since $\det(hy, hv_2, \ldots, hv_n) = \det(h)\det(y, v_2, \ldots, v_n)$. Write $|S^{n-1}| = \int_{S^{n-1}}\sigma$ for the area of the unit sphere.

\begin{lemma}\label{characteristic-classes:lemma-sphere-area}
$|S^{n-1}| = n\,|B^n|$, where $|B^n|$ is the volume of the unit ball, and $|B^{2m}| = \pi^m/m!$. Hence $|S^{2m-1}| = 2\pi^m/(m-1)!$; for instance $|S^1| = 2\pi$ and $|S^3| = 2\pi^2$.
\end{lemma}

\begin{proof}
On $\RR^n$, $d\sigma = \sum_i(-1)^{i-1}dy^i \wedge dy^1 \wedge \cdots \wedge \widehat{dy^i} \wedge \cdots \wedge dy^n = n\,dy^1 \wedge \cdots \wedge dy^n$, so Stokes' theorem
(Theorem~\ref{differential-forms:theorem-stokes}) on the unit ball gives $\int_{S^{n-1}}\sigma = n|B^n|$. For the volumes, slice $B^{n+2} \subseteq \RR^2 \times \RR^n$ by the first two
coordinates: for $w \in \RR^2$ with $|w| < 1$ the slice is a ball of radius $(1 - |w|^2)^{1/2}$ in $\RR^n$, of volume $|B^n|(1 - |w|^2)^{n/2}$ (a linear change of variables,
Lemma~\ref{measure-theory:lemma-linear-change}). By Fubini's theorem (Theorem~\ref{measure-theory:theorem-fubini}) and polar coordinates in the plane
(Example~\ref{measure-theory:example-gaussian}),
\[
|B^{n+2}| = |B^n|\int_0^{2\pi}\!\!\int_0^1(1 - r^2)^{n/2}r\,dr\,d\varphi = |B^n|\,\frac{2\pi}{n + 2} .
\]
Starting from $|B^2| = \pi$, induction gives $|B^{2m}| = \pi^m/m!$, and then $|S^{2m-1}| = 2m\,\pi^m/m! = 2\pi^m/(m-1)!$.
\end{proof}

Now let $M$ be oriented of dimension $n$, equal to the rank of $E$, and let $s$ be a section of $E$ with an isolated zero $z$. Choose a positively oriented chart $\varphi \colon V \to \RR^n$ with
$\varphi(z) = 0$ and a positively oriented frame $e$ of $E$ on $V$, not necessarily orthonormal, and write $s = \sum_is^ie_i$. For small $\varepsilon > 0$ let
$\bar B_\varepsilon = \varphi^{-1}(\{|u| \leq \varepsilon\})$, let $B_\varepsilon$ be its interior, assume $z$ is the only zero of $s$ in $\bar B_\varepsilon$, and let $S_\varepsilon = \partial\bar B_\varepsilon$, oriented as the boundary of
$\bar B_\varepsilon$. On $V \setminus \{z\}$ define $f = (s^1, \ldots, s^n)/|(s^1, \ldots, s^n)|$, with values in $S^{n-1}$.

\begin{definition}\label{characteristic-classes:definition-index}
The \emph{index} of $s$ at the isolated zero $z$ is
\[
\operatorname{ind}_z(s) = \frac{1}{|S^{n-1}|}\int_{S_\varepsilon}f^*\sigma .
\]
For $E = TM$ this is the index of a vector field at an isolated zero.
\end{definition}

\begin{lemma}\label{characteristic-classes:lemma-index}
\begin{enumerate}
\item The index does not depend on $\varepsilon$ nor on the positively oriented frame $e$.
\item If $V \subseteq \RR^n$, $E = V \times \RR^n$ with the standard frame and $s(u) = hu$ near $z = 0$ for some $h \in \mathrm{O}(n)$, then $\operatorname{ind}_0(s) = \det h$. In particular a
source, $s(u) = u$, has index $1$, and a sink, $s(u) = -u$, has index $(-1)^n$.
\end{enumerate}
\end{lemma}

\begin{proof}
Two facts are used. First, $f^*\sigma$ is closed on $V \setminus \{z\}$, being the pullback of a form of top degree on $S^{n-1}$. Second, if $f_t \colon S_\varepsilon \to S^{n-1}$, $t \in [0, 1]$,
is a smooth homotopy, then $\int_{S_\varepsilon}f_0^*\sigma = \int_{S_\varepsilon}f_1^*\sigma$: by Theorem~\ref{de-rham:theorem-homotopy-invariance} the closed forms $f_0^*\sigma$ and $f_1^*\sigma$ differ by an
exact form, whose integral over the compact manifold $S_\varepsilon$ without boundary vanishes by Stokes' theorem.

(1) For $\varepsilon^{\prime} < \varepsilon$, the region $\varepsilon^{\prime} \leq |\varphi| \leq \varepsilon$ is a compact manifold with boundary $S_\varepsilon$ and $S_{\varepsilon^{\prime}}$ with the opposite orientation, so
Stokes' theorem and $d(f^*\sigma) = 0$ give the same integral over both spheres. If $e^{\prime} = eg$ is another positively oriented frame, with $g$ taking values in invertible matrices of
positive determinant, then the components are related by $(s^i) = g\,(s^{\prime i})$. Suppose first that both frames are orthonormal, so that $g$ takes values in $\mathrm{SO}(n)$. Then
$f^{\prime} = g^Tf$, and $f_t(q) = g(\varphi^{-1}(t\varphi(q)))^Tf(q)$, $t \in [0, 1]$, is a smooth homotopy on $S_\varepsilon$ from $f_0 = g(z)^T \circ f$ to $f_1 = f^{\prime}$. Since
$g(z)^T \in \mathrm{SO}(n)$, $f_0^*\sigma = f^*(g(z)^{T*}\sigma) = f^*\sigma$, and the claim follows. In general, apply the Gram--Schmidt process to $e$: it gives an oriented orthonormal frame
$e^{\prime} = eg$ with $g$ upper triangular with positive diagonal entries, so $(s^i) = g\,(s^{\prime i})$. The matrices $g_t = (1 - t)I + tg$ are again upper triangular with positive
diagonal, hence invertible, and $f_t = g_t(s^{\prime i})/|g_t(s^{\prime i})|$ is a homotopy from the map $f^{\prime}$ computed in $e^{\prime}$ to the map $f$ computed in $e$.

(2) Here $f(u) = hu/|u|$, and on $S_\varepsilon$ this is $h$ composed with the dilation $u \mapsto u/\varepsilon$ onto $S^{n-1}$, which preserves orientations. So
$\int_{S_\varepsilon}f^*\sigma = \int_{S^{n-1}}h^*\sigma = \det(h)|S^{n-1}|$. The sink is $h = -I$, of determinant $(-1)^n$.
\end{proof}

The index is the degree of $f \colon S_\varepsilon \to S^{n-1}$ (Exercise~\ref{de-rham:exercise-degree-de-rham}), so it is an integer, and it does not depend on the chart either
\cite[Chapter 6]{milnor-differential-viewpoint}; neither fact is needed below, where any chart may be used at each zero. For a vector field in the plane, $f$ is the direction of the field
along a small circle, and the index is the number of turns it makes: for $X(z) = z^k$ (in complex notation) it is $k$, and for $X(z) = \bar z^k$ it is $-k$
(Exercise~\ref{characteristic-classes:exercise-indices-plane}).

We now come to Chern's construction, which takes place on the unit sphere bundle of $E$.

The \emph{unit sphere bundle} of $E$ is $SE = \{v \in E : |v| = 1\}$, with projection $\pi \colon SE \to M$. Over an open set $U$ with an oriented orthonormal frame $e$, the map
$U \times S^{n-1} \to SE|_U$, $(x, y) \mapsto \sum_iy^ie_i(x)$, is a bijection; for another such frame $e^{\prime} = eg$ the transition map is $(x, y) \mapsto (x, g(x)^Ty)$, a diffeomorphism. So
$SE$ is a manifold of dimension $\dim M + n - 1$ (a closed submanifold of $E$), $\pi$ is a submersion, and the fibre $S_xE$ is the unit sphere of $E_x$. On $SE|_U$ we have the functions
$y^1, \ldots, y^n$, with $\sum_i(y^i)^2 = 1$, and we define the $1$-forms
\[
\eta^i = dy^i + \sum_j\theta_{ij}\,y^j, \qquad i = 1, \ldots, n,
\]
where, here and below, forms on $M$ are pulled back to $SE$ by $\pi$ without changing the notation. The meaning of $\eta$ is the following. On $SE$ the bundle $\pi^*E$ has the
\emph{tautological section} $v \mapsto v$, which is $\sum_iy^i\,e_i$ in the frame $\pi^*e$, and $\eta^i$ are the components of its covariant derivative for the pullback connection
(Proposition~\ref{connections:proposition-connections-exist}(2)): $\eta$ measures how the point $v$ of the sphere bundle moves compared with parallel transport.

\begin{lemma}\label{characteristic-classes:lemma-tautological}
\begin{enumerate}
\item $\sum_iy^i\eta^i = 0$; in particular $\eta^1 \wedge \cdots \wedge \eta^n = 0$.
\item For another oriented orthonormal frame $e^{\prime} = eg$, $y^{\prime} = g^Ty$ and $\eta^{\prime} = g^T\eta$ (as columns, $g$ being composed with $\pi$).
\item $d\eta^i = \sum_jF_{ij}\,y^j - \sum_j\theta_{ij} \wedge \eta^j$.
\item On a fibre $S_xE$, identified with $S^{n-1}$ by the frame at $x$, $\eta^i$ restricts to $dy^i$, and $F$ and $\theta$ restrict to $0$.
\end{enumerate}
\end{lemma}

\begin{proof}
(1) $\sum_iy^i\eta^i = \frac12d\bigl(\sum_i(y^i)^2\bigr) + \sum_{i,j}\theta_{ij}y^iy^j = 0$, the first term because $\sum_i(y^i)^2 = 1$ and the second because $\theta$ is skew-symmetric.
At each point $y \neq 0$, so the $\eta^i$ are linearly dependent and their wedge product vanishes (Proposition~\ref{multilinear-algebra:proposition-exterior-properties}(4)).
(2) From $\sum_iy^ie_i = \sum_jy^{\prime j}e^{\prime}_j = \sum_{i,j}y^{\prime j}g_{ij}e_i$ we get $y = gy^{\prime}$, so $y^{\prime} = g^{-1}y = g^Ty$. Then, in vector notation,
$\eta^{\prime} = dy^{\prime} + \theta^{\prime}y^{\prime} = d(g^{-1})y + g^{-1}dy + (g^{-1}\theta g + g^{-1}dg)g^{-1}y = g^{-1}(dy + \theta y)$, because $d(g^{-1}) = -g^{-1}dg\,g^{-1}$.
(3) $d\eta^i = \sum_jd\theta_{ij}\,y^j - \sum_j\theta_{ij} \wedge dy^j$. Substitute $d\theta_{ij} = F_{ij} - \sum_k\theta_{ik} \wedge \theta_{kj}$ and $dy^j = \eta^j - \sum_k\theta_{jk}y^k$:
the two terms $\mp\sum_{j,k}\theta_{ik} \wedge \theta_{kj}\,y^j$ cancel, leaving the formula.
(4) The composite of the inclusion of $S_xE$ with $\pi$ is constant, so pulled back forms from $M$ of positive degree restrict to zero.
\end{proof}

Now we build forms on $SE$ out of $y$, $\eta$ and $F$, contracting all indices with the sign symbol $\varepsilon_{i_1 \ldots i_n}$, which is the sign of the permutation $(i_1, \ldots, i_n)$ of
$(1, \ldots, n)$, and $0$ if two indices are equal. Let $n = 2m$, and for $0 \leq k \leq m - 1$ put
\[
\Phi_k = \sum\varepsilon_{i_1 \ldots i_n}\,y^{i_1}\,\eta^{i_2} \wedge \cdots \wedge \eta^{i_{n-2k}} \wedge F_{i_{n-2k+1}i_{n-2k+2}} \wedge \cdots \wedge F_{i_{n-1}i_n},
\]
with $n - 2k - 1$ factors $\eta$ and $k$ factors $F$, a form of degree $n - 1$; and for $-1 \leq k \leq m - 1$ put
\[
\Psi_k = \sum\varepsilon_{i_1 \ldots i_n}\,\eta^{i_1} \wedge \cdots \wedge \eta^{i_{n-2k-2}} \wedge F_{i_{n-2k-1}i_{n-2k}} \wedge \cdots \wedge F_{i_{n-1}i_n},
\]
with $n - 2k - 2$ factors $\eta$ and $k + 1$ factors $F$, a form of degree $n$. All sums run over all $(i_1, \ldots, i_n)$.

\begin{lemma}\label{characteristic-classes:lemma-chern-forms}
\begin{enumerate}
\item The forms $\Phi_k$ and $\Psi_k$ do not depend on the oriented orthonormal frame, so they are global forms on $SE$.
\item $\Psi_{-1} = 0$ and $\Psi_{m-1} = 2^mm!\,\pi^*\mathrm{Pf}(F) = (4\pi)^mm!\,\pi^*e(\nabla)$.
\item On a fibre $S_xE \cong S^{n-1}$, $\Phi_k$ restricts to $0$ for $k \geq 1$, and $\Phi_0$ restricts to $(n-1)!\,\sigma$.
\end{enumerate}
\end{lemma}

\begin{proof}
(1) By Lemma~\ref{characteristic-classes:lemma-tautological}(2), in the frame $eg$ we have $y^{\prime i} = \sum_ag_{ai}y^a$, $\eta^{\prime i} = \sum_ag_{ai}\eta^a$ and
$F^{\prime}_{ij} = \sum_{a,b}g_{ai}g_{bj}F_{ab}$. Substituting, every index $i_l$ in $\Phi^{\prime}_k$ becomes a sum over an index $a_l$ with a factor $g_{a_li_l}$, and the order of the
factors is unchanged; so $\Phi^{\prime}_k$ is the same expression as $\Phi_k$ with $\varepsilon_{a_1 \ldots a_n}$ replaced by $\sum_{i_1, \ldots, i_n}\varepsilon_{i_1 \ldots i_n}g_{a_1i_1} \cdots g_{a_ni_n}$.
By the Leibniz formula this is the determinant of the matrix whose rows are the rows $a_1, \ldots, a_n$ of $g$, that is $\varepsilon_{a_1 \ldots a_n}\det g = \varepsilon_{a_1 \ldots a_n}$. The same
argument applies to $\Psi_k$.
(2) $\Psi_{-1} = \sum\varepsilon_{i_1 \ldots i_n}\eta^{i_1} \wedge \cdots \wedge \eta^{i_n} = n!\,\eta^1 \wedge \cdots \wedge \eta^n = 0$ by Lemma~\ref{characteristic-classes:lemma-tautological}(1). And
$\Psi_{m-1} = \sum_\sigma\operatorname{sgn}(\sigma)F_{\sigma(1)\sigma(2)} \wedge \cdots \wedge F_{\sigma(n-1)\sigma(n)} = 2^mm!\,\mathrm{Pf}(F)$.
(3) By Lemma~\ref{characteristic-classes:lemma-tautological}(4) the factors $F$ restrict to $0$ and $\eta^i$ to $dy^i$. So $\Phi_0$ restricts to
$\sum\varepsilon_{i_1 \ldots i_n}y^{i_1}dy^{i_2} \wedge \cdots \wedge dy^{i_n}$. For fixed $i_1 = i$, the $(n-1)!$ orderings of the remaining indices all give
$\varepsilon_{i,1, \ldots, \hat{\imath}, \ldots, n}\,dy^1 \wedge \cdots \wedge \widehat{dy^i} \wedge \cdots \wedge dy^n$, and $\varepsilon_{i,1, \ldots, \hat{\imath}, \ldots, n} = (-1)^{i-1}$; the result is
$(n-1)!\,\sigma$.
\end{proof}

\begin{lemma}[Chern]\label{characteristic-classes:lemma-transgression}
For $0 \leq k \leq m - 1$,
\[
d\Phi_k = \Psi_{k-1} - \frac{2m - 2k - 1}{2k + 2}\,\Psi_k .
\]
\end{lemma}

\begin{proof}
Both sides are global forms on $SE$, so it suffices to compare them at an arbitrary point $v \in SE$, in a frame adapted to $v$. Let $x = \pi(v)$ and choose, by
Lemma~\ref{characteristic-classes:lemma-normal-frames}, an oriented orthonormal frame near $x$ with $\theta(x) = 0$ and $e_n(x) = v$. Then at the point $v$:
\begin{itemize}
\item $y = (0, \ldots, 0, 1)$, and the pulled back forms $\theta_{ij}$ vanish (their value at $v$ is $\theta_{ij}(x) \circ d\pi = 0$);
\item hence $\eta^i = dy^i$, and $\eta^n = dy^n = 0$, because $\sum_iy^idy^i = \frac12d\sum_i(y^i)^2 = 0$ and $y = (0, \ldots, 0, 1)$;
\item $d\eta^i = \sum_jF_{ij}y^j = F_{in}$ by Lemma~\ref{characteristic-classes:lemma-tautological}(3), and $dF = F \wedge \theta - \theta \wedge F = 0$.
\end{itemize}
Write $r = n - 2k$ and $G = F_{i_{r+1}i_{r+2}} \wedge \cdots \wedge F_{i_{n-1}i_n}$ for the product of the $k$ factors $F$ in $\Phi_k$. By the Leibniz rule,
\[
d\bigl(y^{i_1}\eta^{i_2} \wedge \cdots \wedge \eta^{i_r} \wedge G\bigr) = dy^{i_1} \wedge \eta^{i_2} \wedge \cdots \wedge \eta^{i_r} \wedge G
+ y^{i_1}\sum_{l=2}^r(-1)^{l}\,\eta^{i_2} \wedge \cdots \wedge d\eta^{i_l} \wedge \cdots \wedge \eta^{i_r} \wedge G \pm y^{i_1}\eta^{i_2} \wedge \cdots \wedge \eta^{i_r} \wedge dG,
\]
since $d\eta^{i_l}$ is preceded by $l - 2$ factors of degree $1$. We evaluate the three resulting sums at $v$.

\emph{The terms with $dG$} vanish at $v$, since $dF = 0$ there.

\emph{The terms with $dy^{i_1}$.} At $v$, $dy^{i_1} = \eta^{i_1}$, so their sum is $\sum\varepsilon_{i_1 \ldots i_n}\eta^{i_1} \wedge \cdots \wedge \eta^{i_r} \wedge G$, which has $n - 2k$ factors $\eta$ and $k$
factors $F$: it is $\Psi_{k-1}$ at $v$.

\emph{The terms with $d\eta^{i_l}$.} At $v$, $y^{i_1} = \delta_{i_1n}$ and $d\eta^{i_l} = F_{i_ln}$. Moving the $2$-form $F_{i_ln}$ to the right of the remaining factors $\eta$ costs no sign, so
these terms add up to
\[
\sum_{l=2}^r(-1)^l\sum\varepsilon_{n\,i_2 \ldots i_n}\,\eta^{i_2} \wedge \cdots \widehat{\eta^{i_l}} \cdots \wedge \eta^{i_r} \wedge F_{i_ln} \wedge G .
\]
In the $l$-th term rename the indices so that $i_l$ comes last among $i_2, \ldots, i_r$; this moves it past $r - l$ indices and multiplies $\varepsilon$ by $(-1)^{r-l}$. Since $r$ is even,
every $l$ contributes $(-1)^l(-1)^{r-l} = 1$ times the same sum, so the terms with $d\eta$ give $(r-1)\,T$ with
\[
T = \sum\varepsilon_{n\,j_1 \ldots j_{r-2}\,a\,j_{r+1} \ldots j_n}\,\eta^{j_1} \wedge \cdots \wedge \eta^{j_{r-2}} \wedge F_{an} \wedge F_{j_{r+1}j_{r+2}} \wedge \cdots \wedge F_{j_{n-1}j_n} .
\]
Now look at $\Psi_k$ at $v$. Its factors $\eta$ carry the indices in the positions $1, \ldots, r - 2$, and since $\eta^n = 0$ at $v$, the index $n$ must occur in one of the $k + 1$ factors $F$.
Exchanging two pairs of positions among the factors $F$ is an even permutation of the indices and does not change the product of the commuting $2$-forms; exchanging the two positions
of one pair changes the sign of $\varepsilon$ and of that factor $F$. So the terms in which $n$ occupies any of the $2(k+1)$ positions of the factors $F$ add up to the same as those in
which it occupies position $r$, and
\[
\Psi_k = 2(k+1)\sum\varepsilon_{j_1 \ldots j_{r-2}\,a\,n\,j_{r+1} \ldots j_n}\,\eta^{j_1} \wedge \cdots \wedge \eta^{j_{r-2}} \wedge F_{an} \wedge F_{j_{r+1}j_{r+2}} \wedge \cdots \wedge F_{j_{n-1}j_n} \quad\text{at } v.
\]
Moving $n$ from the first position to position $r$ passes $r - 1$ indices, an odd number, so $\varepsilon_{n\,j_1 \ldots j_{r-2}\,a \ldots} = -\varepsilon_{j_1 \ldots j_{r-2}\,a\,n \ldots}$ and
$T = -\Psi_k/(2(k+1))$ at $v$. Altogether, at $v$ we get $d\Phi_k = \Psi_{k-1} - \frac{r - 1}{2(k+1)}\Psi_k$, and $r - 1 = 2m - 2k - 1$.
\end{proof}

Write $(2m-1)!! = 1 \cdot 3 \cdot 5 \cdots (2m-1)$.

\begin{proposition}[Chern]\label{characteristic-classes:proposition-transgression}
Let $n = 2m$, $c_k = \frac{2m-2k-1}{2k+2}$, $a_0 = 1$ and $a_k = c_0c_1 \cdots c_{k-1}$, and put
\[
\Pi = -\frac{1}{(2\pi)^m(2m-1)!!}\sum_{k=0}^{m-1}a_k\Phi_k \in \Omega^{n-1}(SE) .
\]
Then $d\Pi = \pi^*e(\nabla)$, and the restriction of $\Pi$ to every fibre $S_xE \cong S^{n-1}$ is $-\sigma/|S^{n-1}|$; in particular $\int_{S_xE}\Pi = -1$.
\end{proposition}

\begin{proof}
By Lemma~\ref{characteristic-classes:lemma-transgression},
\[
d\sum_{k=0}^{m-1}a_k\Phi_k = \sum_{k=0}^{m-1}a_k\Psi_{k-1} - \sum_{k=0}^{m-1}a_kc_k\Psi_k = \sum_{k=0}^{m-2}(a_{k+1} - a_kc_k)\Psi_k - a_{m-1}c_{m-1}\Psi_{m-1} = -a_m\Psi_{m-1},
\]
using $\Psi_{-1} = 0$ and $a_{k+1} = a_kc_k$. Here
$a_m = \prod_{k=0}^{m-1}\frac{2m-2k-1}{2k+2} = \frac{(2m-1)(2m-3) \cdots 1}{2^mm!} = \frac{(2m-1)!!}{2^mm!}$, and $\Psi_{m-1} = 2^mm!(2\pi)^m\pi^*e(\nabla)$ by
Lemma~\ref{characteristic-classes:lemma-chern-forms}(2). So $d\sum_ka_k\Phi_k = -(2m-1)!!(2\pi)^m\pi^*e(\nabla)$, which is $d\Pi = \pi^*e(\nabla)$. On a fibre only $\Phi_0$ survives, with
$a_0 = 1$, by Lemma~\ref{characteristic-classes:lemma-chern-forms}(3), so $\Pi$ restricts to $-\frac{(2m-1)!}{(2\pi)^m(2m-1)!!}\sigma$. Now $(2m-1)!/(2m-1)!! = 2 \cdot 4 \cdots (2m-2) = 2^{m-1}(m-1)!$,
and by Lemma~\ref{characteristic-classes:lemma-sphere-area},
$2^{m-1}(m-1)!/(2\pi)^m = (m-1)!/(2\pi^m) = 1/|S^{2m-1}|$.
\end{proof}

So the pullback of the Euler form to the sphere bundle is exact, although the Euler form itself is in general not exact on $M$; a form $\Pi$ with these two properties is called a
\emph{global angular form} \cite[Section 11]{bott-tu}, because on each fibre it is (minus) the normalised angular form of the sphere. As a consequence, $\pi^*d\,e(\nabla) = d\,d\Pi = 0$, and
$\pi^*$ is injective on forms because $\pi$ is a submersion; this gives a second proof that $e(\nabla)$ is closed.

\begin{example}\label{characteristic-classes:example-chern-surfaces}
Let $n = 2$, so $m = 1$, $\Pi = -\frac{1}{2\pi}\Phi_0$ and $\Phi_0 = y^1\eta^2 - y^2\eta^1$. Over a set $U$ with an oriented orthonormal frame, a point of $SE|_U$ is
$v = \cos\phi\,e_1 + \sin\phi\,e_2$, where the angle $\phi$ is a local function on $SE|_U$ defined up to adding constant multiples of $2\pi$, so $d\phi$ is well defined on $SE|_U$.
Then $y = (\cos\phi, \sin\phi)$, and with $\theta_{21} = -\theta_{12}$,
\[
\eta^1 = -\sin\phi\,d\phi + \theta_{12}\sin\phi, \qquad \eta^2 = \cos\phi\,d\phi + \theta_{21}\cos\phi, \qquad \Phi_0 = d\phi + \theta_{21} .
\]
So $\Pi = -\frac{1}{2\pi}(d\phi + \theta_{21})$ and $d\Pi = -\frac{1}{2\pi}d\theta_{21} = \frac{1}{2\pi}F_{12} = \pi^*e(\nabla)$, as $F_{12} = d\theta_{12}$ for $n = 2$. The form $d\phi + \theta_{21}$
measures the turning of the unit vector $v$ relative to parallel transport: along a curve, $v$ is parallel exactly when $\dot\phi = -\theta_{21}(\dot\gamma)$, since
$\nabla e_1 = \theta_{21} \otimes e_2$.

The proof of Theorem~\ref{characteristic-classes:theorem-chern} below becomes, in this case, the following classical argument. Let $s$ be a section with isolated zeros, and on
$M \setminus \{s = 0\}$ use the frame $e_1 = s/|s|$, $e_2$ its rotation by $+\pi/2$. Then $s/|s|$ has angle $\phi = 0$ in this frame, so
$\int_{M \setminus \cup B_\varepsilon}e(\nabla) = \frac{1}{2\pi}\sum_z\int_{S_\varepsilon(z)}\theta_{21}$, by Stokes' theorem and $\frac{1}{2\pi}F_{12} = -\frac{1}{2\pi}d\theta_{21}$. Near a zero $z$ compare with a
frame $(e^{\prime}_1, e^{\prime}_2)$ defined on all of $B_\varepsilon(z)$: if $e_1 = \cos\psi\,e^{\prime}_1 + \sin\psi\,e^{\prime}_2$, the connection forms are related by $\theta_{21} = \theta^{\prime}_{21} + d\psi$
(Exercise~\ref{characteristic-classes:exercise-rotation-gauge}). The integral of $d\psi$ over $S_\varepsilon(z)$ is $2\pi\operatorname{ind}_z(s)$, since $\psi$ is the angle of $s$ in the frame
$e^{\prime}$, and the integral of the smooth form $\theta^{\prime}_{21}$ tends to $0$ with $\varepsilon$. So $\int_Me(\nabla) = \sum_z\operatorname{ind}_z(s)$.
\end{example}

\begin{theorem}[Chern]\label{characteristic-classes:theorem-chern}
Let $M$ be a compact oriented manifold of even dimension $n = 2m$, $E \to M$ an oriented real vector bundle of rank $n$ with a metric and a metric connection $\nabla$, and $s$ a section of
$E$ whose zeros are isolated. Then $s$ has finitely many zeros, and
\[
\int_Me(\nabla) = \sum_{s(z) = 0}\operatorname{ind}_z(s) .
\]
\end{theorem}

\begin{proof}
The zero set $Z$ of $s$ is closed in the compact space $M$ and discrete, hence finite. For each $z \in Z$ choose a positively oriented chart $\varphi_z$ centred at $z$ and an oriented
orthonormal frame $e$ of $E$ on its domain, and for small $\varepsilon > 0$ let $\bar B_\varepsilon(z)$ and $S_\varepsilon(z)$ be as in Definition~\ref{characteristic-classes:definition-index},
with the balls pairwise disjoint. The section $\hat s = s/|s|$ of $SE$ over $M \setminus Z$ is a smooth map $\hat s \colon M \setminus Z \to SE$ with $\pi \circ \hat s = \id$, so
$\hat s^*\pi^*e(\nabla) = e(\nabla)$. Let $M_\varepsilon = M \setminus \bigcup_zB_\varepsilon(z)$, a compact manifold with boundary the union of the spheres $S_\varepsilon(z)$, each with the
orientation opposite to its orientation as the boundary of $\bar B_\varepsilon(z)$. By Proposition~\ref{characteristic-classes:proposition-transgression} and Stokes' theorem,
\[
\int_{M_\varepsilon}e(\nabla) = \int_{M_\varepsilon}\hat s^*d\Pi = \int_{M_\varepsilon}d(\hat s^*\Pi) = -\sum_{z \in Z}\int_{S_\varepsilon(z)}\hat s^*\Pi .
\]
As $\varepsilon \to 0$ the left side tends to $\int_Me(\nabla)$, since $e(\nabla)$ is a continuous form and the volumes of the balls in the charts tend to $0$. It remains to show that
$\int_{S_\varepsilon(z)}\hat s^*\Pi \to -\operatorname{ind}_z(s)$ for each zero $z$.

Fix $z$, write $S_\varepsilon = S_\varepsilon(z)$, and use the frame $e$ to identify $SE$ over the chart domain with the product of the domain and $S^{n-1}$. In this identification
$\hat s(q) = (q, f(q))$, with $f$ as in Definition~\ref{characteristic-classes:definition-index}. Define
\[
H \colon S_\varepsilon \times [0, 1] \to SE, \qquad H(q, t) = \bigl(\varphi_z^{-1}(t\varphi_z(q)),\, f(q)\bigr),
\]
so that $H(q, 1) = \hat s(q)$ and $H(q, 0) = \iota_z(f(q))$, where $\iota_z \colon S^{n-1} \to S_zE$ is the identification of the fibre given by the frame at $z$. By Stokes' theorem on the
compact manifold $S_\varepsilon \times [0, 1]$, whose boundary is $S_\varepsilon \times \{1\}$ and $S_\varepsilon \times \{0\}$ with opposite orientations,
\[
\int_{S_\varepsilon}\hat s^*\Pi - \int_{S_\varepsilon}f^*\iota_z^*\Pi = \pm\int_{S_\varepsilon \times [0, 1]}H^*d\Pi = \pm\int_{S_\varepsilon \times [0, 1]}(\pi \circ H)^*e(\nabla) .
\]
The map $\pi \circ H(q, t) = \varphi_z^{-1}(t\varphi_z(q))$ restricts to a diffeomorphism from $S_\varepsilon \times (0, 1]$ onto $\bar B_\varepsilon(z) \setminus \{z\}$ (in the chart it is
$(u, t) \mapsto tu$, with inverse $w \mapsto (\varepsilon w/|w|, |w|/\varepsilon)$), and $S_\varepsilon \times \{0\}$ has measure zero; so the right side is $\pm\int_{B_\varepsilon(z)}e(\nabla)$,
which tends to $0$ with $\varepsilon$. On the other hand, by Proposition~\ref{characteristic-classes:proposition-transgression}, $\iota_z^*\Pi = -\sigma/|S^{n-1}|$, so
$\int_{S_\varepsilon}f^*\iota_z^*\Pi = -\frac{1}{|S^{n-1}|}\int_{S_\varepsilon}f^*\sigma = -\operatorname{ind}_z(s)$ by Lemma~\ref{characteristic-classes:lemma-index}(1). Hence
$\int_{S_\varepsilon}\hat s^*\Pi \to -\operatorname{ind}_z(s)$, and $\int_Me(\nabla) = \sum_z\operatorname{ind}_z(s)$.
\end{proof}

The theorem has a striking structure: the left side does not depend on the section, and the right side does not depend on the metric or the connection. So both are invariants of
the oriented bundle $E$ alone; the left side shows it is the integral of a characteristic form, the right side that it is an integer. A section without zeros forces it to vanish, in
accordance with Proposition~\ref{characteristic-classes:proposition-euler-de-rham}(3). The de Rham class of $e(\nabla)$ is the Euler class of
Definition~\ref{characteristic-classes:definition-euler-de-rham} \cite[Appendix C]{milnor-stasheff}; we do not need this.

\begin{corollary}[Gauss--Bonnet--Chern]\label{characteristic-classes:corollary-gauss-bonnet-chern}
Let $M$ be a compact oriented manifold of dimension $2m$ with a Riemannian metric, and $\nabla$ a metric connection on $TM$, for instance the Levi-Civita connection.
\begin{enumerate}
\item For every vector field $X$ on $M$ with isolated zeros, $\int_Me(TM, \nabla) = \sum_z\operatorname{ind}_z(X)$.
\item (Hopf) The sum of the indices of a vector field with isolated zeros is the same for all such vector fields.
\item $\int_Me(TM, \nabla) = \chi(M)$. In particular the integral does not depend on the metric, and if $M$ admits a nowhere vanishing vector field, then $\chi(M) = 0$.
\end{enumerate}
\end{corollary}

\begin{proof}
(1) is Theorem~\ref{characteristic-classes:theorem-chern} for $E = TM$, and (2) follows because the left side of (1) does not involve $X$. For (3), the sum of the indices equals $\chi(M)$ by the
Poincar\'e--Hopf theorem, Theorem~\ref{characteristic-classes:theorem-euler}(3), applied to any vector field with isolated zeros (for instance the gradient of a Morse function
\cite{milnor-morse}). A nowhere vanishing vector field has no zeros, so the sum of its indices is $0$.
\end{proof}

For surfaces, $e(TM, \nabla) = \frac{1}{2\pi}F_{12}$, and for the Levi-Civita connection $F_{12} = K\,dA$ with $K$ the Gaussian curvature: this is the Gauss--Bonnet theorem
$\int_MK\,dA = 2\pi\chi(M)$, proved again in Chapter~\ref{riemannian-geometry} together with its higher-dimensional Riemannian form. There the curvature of the round sphere $S^{2m}$ is
computed, and Corollary~\ref{characteristic-classes:corollary-gauss-bonnet-chern} gives a formula for its volume.

\section{Stiefel--Whitney classes and Wu's formula}\label{characteristic-classes:section-wu}

With coefficients in $\mathbb{F}_2$ no orientations are needed, and the Steenrod squares
(Theorem~\ref{cohomology-operations:theorem-steenrod-squares}) produce the Stiefel--Whitney classes directly from the Thom class. For the tangent
bundle of a closed manifold this has a striking consequence, Wu's formula: the Stiefel--Whitney classes of $TM$ are determined by the cohomology
ring $H^*(M; \mathbb{F}_2)$ with its Steenrod squares, so they are homotopy invariants, although $TM$ is not. In this section $H^*$ means
$H^*(-; \mathbb{F}_2)$, and $\mathrm{Sq} = \sum_i\mathrm{Sq}^i$.

\begin{theorem}[Thom]\label{characteristic-classes:theorem-stiefel-whitney-thom}
Let $\pi \colon E \to M$ be a real vector bundle of rank $n$ over a manifold, and $E_0$ the complement of the zero section.
\begin{enumerate}
\item (Thom isomorphism mod $2$) There is a unique class $u \in H^n(E, E_0)$ restricting to the nonzero class on every fibre, and
$\Phi(x) = \pi^*(x) \smile u$ is an isomorphism $H^k(M) \to H^{k+n}(E, E_0)$ for all $k$.
\item The classes $w_i(E) = \Phi^{-1}(\mathrm{Sq}^i u) \in H^i(M)$ are the Stiefel--Whitney classes of
Definition~\ref{characteristic-classes:definition-pontryagin}: $w_0 = 1$, $w_i = 0$ for $i > n$, they are natural, the total class
$w = \sum_iw_i$ satisfies $w(E \oplus E') = w(E)w(E')$, and $w_1$ of the tautological line bundle on $\RR\PP^1$ is nonzero. Equivalently,
$\mathrm{Sq}(u) = \pi^*w(E) \smile u$.
\end{enumerate}
\end{theorem}

This is \cite[Chapters 8 and 10]{milnor-stasheff}; it is Thom's definition of the Stiefel--Whitney classes. The properties in (2) follow from the
axioms of Theorem~\ref{cohomology-operations:theorem-steenrod-squares}: $w_0 = 1$ from $\mathrm{Sq}^0 = \id$, $w_i = 0$ for $i > n$ from
$\mathrm{Sq}^iu = 0$ for $i > \deg u$, and the Whitney formula from the Cartan formula, as the Thom class of $E \oplus E'$ is the cross product of
the Thom classes.

\begin{definition}\label{characteristic-classes:definition-wu-classes}
Let $M$ be a closed $n$-manifold and $[M] \in H_n(M)$ its fundamental class (every manifold is $\mathbb{F}_2$-oriented,
Proposition~\ref{poincare-duality:proposition-orientation-covering}(3)). For each $i$, the map $H^{n-i}(M) \to \mathbb{F}_2$,
$x \mapsto \langle\mathrm{Sq}^i(x), [M]\rangle$, is linear, so by Poincar\'e duality (Corollary~\ref{poincare-duality:corollary-cup-pairing}) there is a
unique class $v_i \in H^i(M)$ with
\[
\langle v_i \smile x, [M]\rangle = \langle\mathrm{Sq}^i(x), [M]\rangle \quad \text{for all } x \in H^{n-i}(M).
\]
The $v_i$ are the \emph{Wu classes} of $M$, and $v = 1 + v_1 + v_2 + \cdots$ is the \emph{total Wu class}.
\end{definition}

Since $\mathrm{Sq}^0 = \id$, $v_0 = 1$. If $2i > n$, then $\mathrm{Sq}^i(x) = 0$ for $x$ of degree $n - i < i$, so $v_i = 0$: only
$v_0, \ldots, v_{\lfloor n/2\rfloor}$ can be nonzero.

\begin{theorem}[Wu's formula]\label{characteristic-classes:theorem-wu}
Let $M$ be a closed smooth $n$-manifold. Then $w(TM) = \mathrm{Sq}(v)$, that is, $w_k(TM) = \sum_i\mathrm{Sq}^{k-i}(v_i)$ for every $k$.
\end{theorem}

\begin{proof}[Sketch of proof]
Choose a basis $(b_j)$ of $H^*(M)$ and the dual basis $(b_j^\#)$ for the cup pairing, $\langle b_j \smile b_l^\#, [M]\rangle = \delta_{jl}$
(Corollary~\ref{poincare-duality:corollary-cup-pairing}); then every $y \in H^*(M)$ is $y = \sum_j\langle y \smile b_j, [M]\rangle\,b_j^\#$. We quote two
geometric facts \cite[Chapter 11]{milnor-stasheff}, which are the omitted verifications. First, the class $U = \sum_jb_j \times b_j^\# \in H^n(M \times M)$
(K\"unneth, Theorem~\ref{singular-cohomology:theorem-kunneth}) is the Poincar\'e dual of the diagonal. Second, a tubular neighbourhood of the diagonal is
isomorphic to $TM$, and under this isomorphism and excision $U$ corresponds to the Thom class of $TM$; with
Theorem~\ref{characteristic-classes:theorem-stiefel-whitney-thom}(2) this gives $\mathrm{Sq}(U) = (1 \times w(TM)) \smile U$.

The rest is algebra. Let $\Theta \colon H^*(M \times M) \to H^*(M)$ be the linear map with $\Theta(a \times b) = \langle a, [M]\rangle\,b$ (the slant
product with $[M]$ in the first factor), so that $\Theta((1 \times w) \smile Z) = w \smile \Theta(Z)$. Since $\langle b_j, [M]\rangle = \langle 1 \smile b_j, [M]\rangle$,
$\Theta(U) = \sum_j\langle 1 \smile b_j, [M]\rangle b_j^\# = 1$, so $\Theta(\mathrm{Sq}(U)) = w(TM)$. On the other hand, by the Cartan formula
$\mathrm{Sq}(U) = \sum_j\mathrm{Sq}(b_j) \times \mathrm{Sq}(b_j^\#)$, and $\langle\mathrm{Sq}(b_j), [M]\rangle = \langle\mathrm{Sq}^{n - \deg b_j}(b_j), [M]\rangle
= \langle v \smile b_j, [M]\rangle$ by the definition of the Wu classes; so
$\Theta(\mathrm{Sq}(U)) = \sum_j\langle v \smile b_j, [M]\rangle\,\mathrm{Sq}(b_j^\#) = \mathrm{Sq}\bigl(\sum_j\langle v \smile b_j, [M]\rangle b_j^\#\bigr) = \mathrm{Sq}(v)$.
\end{proof}

\begin{corollary}\label{characteristic-classes:corollary-wu-invariance}
Let $M$ and $N$ be closed connected smooth $n$-manifolds and $h \colon M \to N$ a homotopy equivalence. Then $h^*w(TN) = w(TM)$, and $M$ and $N$
have the same Stiefel--Whitney numbers $\langle w_{i_1}(T-) \cdots w_{i_r}(T-), [-]\rangle$. Moreover $w_1(TM) = v_1$, so $M$ is orientable if and only if
$\mathrm{Sq}^1 \colon H^{n-1}(M) \to H^n(M)$ is zero.
\end{corollary}

\begin{proof}
The map $h^*$ is an isomorphism of rings commuting with the Steenrod squares (naturality), and $h_*[M] = [N]$, because $h_*$ is an isomorphism
$H_n(M) \cong \mathbb{F}_2 \to H_n(N) \cong \mathbb{F}_2$ (Theorem~\ref{poincare-duality:theorem-fundamental-class}). For $x \in H^{n-i}(N)$,
$\langle h^*(v_i(N)) \smile h^*x, [M]\rangle = \langle v_i(N) \smile x, [N]\rangle = \langle\mathrm{Sq}^i(x), [N]\rangle = \langle\mathrm{Sq}^i(h^*x), [M]\rangle$, and $h^*$ is
surjective; so $h^*v(N) = v(M)$, and $h^*w(TN) = h^*\mathrm{Sq}(v(N)) = \mathrm{Sq}(v(M)) = w(TM)$ by Theorem~\ref{characteristic-classes:theorem-wu}. The
Stiefel--Whitney numbers agree because $\langle h^*y, [M]\rangle = \langle y, h_*[M]\rangle$. Next, $w_1 = \mathrm{Sq}^1(v_0) + \mathrm{Sq}^0(v_1) = v_1$, as
$\mathrm{Sq}^1(1) = 0$. By definition $v_1 = 0$ if and only if $\langle\mathrm{Sq}^1(x), [M]\rangle = 0$ for all $x \in H^{n-1}(M)$, that is, since
$\langle -, [M]\rangle$ is an isomorphism $H^n(M) \to \mathbb{F}_2$, if and only if $\mathrm{Sq}^1 = 0$ on $H^{n-1}(M)$; and $w_1(TM) = 0$ if and only if $TM$,
equivalently $M$, is orientable (Definition~\ref{characteristic-classes:definition-pontryagin}).
\end{proof}

\begin{lemma}\label{characteristic-classes:lemma-characteristic-element}
Let $V$ be a finite-dimensional vector space over $\mathbb{F}_2$ with a nondegenerate symmetric bilinear form $B$. There is a unique $c \in V$ with
$B(x, x) = B(c, x)$ for all $x$, and $B(c, c) \equiv \dim V \pmod 2$.
\end{lemma}

\begin{proof}
The map $x \mapsto B(x, x)$ is linear, because $B(x + y, x + y) = B(x, x) + B(y, y) + 2B(x, y)$, so $c$ exists and is unique by nondegeneracy. We
argue by induction on $\dim V$, the case $V = 0$ being trivial. If $B(x, x) = 0$ for all $x$, then $c = 0$ and $B$ is alternating. Choose $x \neq 0$ and
$y$ with $B(x, y) = 1$; the plane $P$ spanned by $x$ and $y$ has Gram matrix $\bigl(\begin{smallmatrix}0 & 1\\ 1 & 0\end{smallmatrix}\bigr)$, so $B$ is
nondegenerate on $P$, $V = P \oplus P^{\perp_B}$ (Lemma~\ref{linear-algebra:lemma-orthogonal-complement}), $B$ is nondegenerate and alternating on
$P^{\perp_B}$, and by induction $\dim V$ is even, as required. Otherwise choose $x$ with $B(x, x) = 1$; then $V = \mathbb{F}_2x \oplus W$ with
$W = x^{\perp_B}$, $B$ is nondegenerate on $W$, and if $c'$ is the element of $W$ given by induction, $c = x + c'$ works: for $y = \lambda x + w$,
$B(y, y) = \lambda + B(w, w) = \lambda + B(c', w) = B(x + c', y)$. Then $B(c, c) = 1 + B(c', c') \equiv 1 + \dim W = \dim V$.
\end{proof}

\begin{corollary}\label{characteristic-classes:corollary-top-stiefel-whitney}
For a closed smooth $n$-manifold $M$, $\langle w_n(TM), [M]\rangle \equiv \chi(M) \pmod 2$.
\end{corollary}

\begin{proof}
We may assume $M$ connected. By Theorem~\ref{characteristic-classes:theorem-wu}, $w_n = \sum_i\mathrm{Sq}^{n-i}(v_i)$. For $i < n/2$ the term vanishes,
because $n - i > i = \deg v_i$, and for $i > n/2$ we have $v_i = 0$. If $n$ is odd nothing remains, so $w_n = 0$, and $\chi(M) = 0$ by
Corollary~\ref{poincare-duality:corollary-euler-odd} (closed manifolds have finitely generated homology \cite[Appendix]{hatcher-at}). If $n = 2k$, then
$w_n = \mathrm{Sq}^k(v_k) = v_k^2$. The cup pairing $B(x, y) = \langle x \smile y, [M]\rangle$ on $V = H^k(M)$ is symmetric (mod $2$ there are no signs) and
nondegenerate, and $B(x, x) = \langle\mathrm{Sq}^k(x), [M]\rangle = B(v_k, x)$; so $v_k$ is the element $c$ of
Lemma~\ref{characteristic-classes:lemma-characteristic-element}, and $\langle w_n, [M]\rangle = B(v_k, v_k) \equiv \dim H^k(M)$. Finally
$\chi(M) = \sum_i(-1)^i\dim H^i(M)$ (Corollary~\ref{poincare-duality:corollary-euler-odd} and universal coefficients over a field), and
$\dim H^i(M) = \dim H^{n-i}(M)$ by Poincar\'e duality; so the terms with $i \neq k$ cancel in pairs modulo $2$, and $\chi(M) \equiv \dim H^k(M)$.
\end{proof}

\begin{example}\label{characteristic-classes:example-wu-projective}
Let $M = \RR\PP^n$, with $H^*(M) = \mathbb{F}_2[\alpha]/(\alpha^{n+1})$, $[M]$ pairing to $1$ with $\alpha^n$, and $\mathrm{Sq}^i(\alpha^m) = \binom{m}{i}\alpha^{m+i}$
(Proposition~\ref{cohomology-operations:proposition-squares-projective}, restricted from $\RR\PP^\infty$). Then $v_i = \binom{n-i}{i}\alpha^i$, since
$\mathrm{Sq}^i(\alpha^{n-i}) = \binom{n-i}{i}\alpha^n$, and $\mathrm{Sq}(\alpha^i) = \alpha^i(1 + \alpha)^i$.
\begin{enumerate}
\item $n = 2$: $v = 1 + \alpha$ and $w = 1 + \alpha(1 + \alpha) = 1 + \alpha + \alpha^2$. So $\RR\PP^2$ is not orientable ($w_1 \neq 0$), and
$w_2 \neq 0$ matches $\chi(\RR\PP^2) = 1$.
\item $n = 3$: $v_1 = \binom{2}{1}\alpha = 0$, so $v = 1$ and $w = 1$: all Stiefel--Whitney classes of $\RR\PP^3$ vanish. This is consistent with
$\RR\PP^3 \cong \mathrm{SO}(3)$ being a Lie group, hence parallelisable.
\item $n = 4$: $v = 1 + \alpha + \alpha^2$ and $w = 1 + \alpha(1 + \alpha) + \alpha^2(1 + \alpha)^2 = 1 + \alpha + \alpha^4$. So $w_4[\RR\PP^4] = 1$, and
$\RR\PP^4$ has a nonzero Stiefel--Whitney number; by a theorem of Pontryagin, the Stiefel--Whitney numbers of a boundary vanish, so
$\RR\PP^4$ is not the boundary of any compact $5$-manifold \cite[Chapter 4]{milnor-stasheff}. (Thom proved the converse: a closed manifold whose
Stiefel--Whitney numbers all vanish is a boundary \cite{thom-1954}.)
\end{enumerate}
In all three cases $w = (1 + \alpha)^{n+1}$, which is the general formula for $\RR\PP^n$ \cite[Chapter 4]{milnor-stasheff}.
\end{example}

\section{Exercises}\label{characteristic-classes:section-exercises}

\begin{exercise}\label{characteristic-classes:exercise-chern-projective}
Using the Euler sequence, compute $c(T\CC\PP^2)$ and $\int_{\CC\PP^2}c_2(T\CC\PP^2)$, and check that the latter is $\chi(\CC\PP^2) = 3$.
\end{exercise}

\begin{solution}
$c(T\CC\PP^2) = (1 + h)^3 = 1 + 3h + 3h^2$, so $c_2 = 3h^2$ and $\int_{\CC\PP^2}c_2 = 3$ since $\int h^2 = 1$. By Theorem~\ref{characteristic-classes:theorem-euler}, $c_2(T\CC\PP^2) = e(T\CC\PP^2)$ integrates to
$\chi(\CC\PP^2) = 3$ (Example~\ref{singular-homology:example-euler}).
\end{solution}

\begin{exercise}\label{characteristic-classes:exercise-flat-bundle-classes}
Show that a bundle with a flat connection has vanishing Chern classes in de Rham cohomology, and give an example where the integral Chern classes are nonzero.
\end{exercise}

\begin{solution}
Chern--Weil forms are computed from the curvature, which is zero. A flat line bundle on a space with torsion in $H^2(M; \ZZ)$, for instance a nontrivial real line bundle on $\RR\PP^2$ complexified, has
torsion $c_1 \neq 0$ mapping to zero in de Rham cohomology.
\end{solution}

\begin{exercise}\label{characteristic-classes:exercise-chern-character-line}
Show that $\mathrm{ch}(L) = e^{c_1(L)}$ for a line bundle, and compute $\mathrm{ch}$ and $\mathrm{td}$ of $T\CC\PP^1$.
\end{exercise}

\begin{exercise}\label{characteristic-classes:exercise-no-almost-complex}
Show that $S^4$ admits no almost complex structure, using that $c_1$ would lie in $H^2(S^4) = 0$ and $c_2 = e(TS^4)$ would have to satisfy $\int_{S^4}c_2 = \chi(S^4) = 2$ while also
$c_1^2 = p_1 + 2c_2$ with $p_1(TS^4) = 0$.
\end{exercise}

\begin{exercise}\label{characteristic-classes:exercise-chern-sum-cp1}
On $\CC\PP^1$ compute the total Chern class of $\mathcal{O}(a) \oplus \mathcal{O}(b)$, and show that $\mathcal{O}(1) \oplus \mathcal{O}(-1)$ and $\mathcal{O} \oplus \mathcal{O}$ have the same Chern classes.
\end{exercise}

\begin{solution}
By the Whitney formula $c = (1 + ah)(1 + bh) = 1 + (a + b)h$, as $h^2 = 0$. For $a = 1$, $b = -1$ this is $1$, the total Chern class of the trivial bundle. (The two bundles are in
fact isomorphic as smooth bundles, but not as holomorphic ones: the first has a holomorphic section vanishing at one point and the second does not.)
\end{solution}

\begin{exercise}\label{characteristic-classes:exercise-c1-determinant}
Show that $c_1(E) = c_1(\det E)$ for every complex vector bundle $E$, using the connection induced on $\det E$ (Proposition~\ref{connections:proposition-induced-connections}(5)).
\end{exercise}

\begin{solution}
If $\nabla$ has connection form $\theta$ in a frame $e$, the induced connection on $\det E$ has connection form $\operatorname{tr}\theta$ in the frame $e_1 \wedge \cdots \wedge e_r$, so its curvature is
$d\operatorname{tr}\theta = \operatorname{tr}(d\theta) = \operatorname{tr}(d\theta + \theta \wedge \theta) = \operatorname{tr}F$, using $\operatorname{tr}(\theta \wedge \theta) = 0$. Hence
$c_1(\det E, \nabla^{\det}) = \frac{i}{2\pi}\operatorname{tr}F = c_1(E, \nabla)$.
\end{solution}

\begin{exercise}\label{characteristic-classes:exercise-pullback-curvature}
Let $f \colon \CC\PP^1 \to \CC\PP^1$ be $z \mapsto z^k$ ($k \geq 1$) and $\nabla$ the projection connection on $\mathcal{O}(-1)$. Show that $\int_{\CC\PP^1}c_1(f^*\mathcal{O}(-1), f^*\nabla) = -k$,
by computing the integral of $f^*c_1(\mathcal{O}(-1), \nabla)$, and compare with Exercise~\ref{vector-bundles:exercise-pullback-degree}.
\end{exercise}

\begin{solution}
The curvature of $f^*\nabla$ is $f^*F$ (Proposition~\ref{connections:proposition-pullback-tensor}(1)), so the Chern form is $f^*\omega$ with $\omega = c_1(\mathcal{O}(-1), \nabla)$. On
$\CC^\times$, $f$ is a $k$-sheeted covering and a local diffeomorphism preserving the orientation (it is holomorphic). Cut $\CC^\times$ into the $k$ open sectors
$\{2\pi j/k < \arg z < 2\pi(j + 1)/k\}$; each is mapped diffeomorphically, preserving orientation, onto $\CC$ minus a ray, and the rays and the points $0$, $\infty$ have measure zero. By
Proposition~\ref{differential-forms:proposition-integral-parametrisation} each sector contributes $\int_{\CC\PP^1}\omega = -1$, so the total is $-k$, in accordance with
$f^*\mathcal{O}(-1) \cong \mathcal{O}(-k)$.
\end{solution}

\begin{exercise}\label{characteristic-classes:exercise-c2-formula}
Verify the formula $c_2 = \frac{1}{8\pi^2}(\operatorname{tr}(F \wedge F) - \operatorname{tr}F \wedge \operatorname{tr}F)$ for a direct sum of two line bundles with curvatures $F_1$, $F_2$, and check that it
equals $c_1(L_1) \wedge c_1(L_2)$.
\end{exercise}

\begin{solution}
Here $F = \mathrm{diag}(F_1, F_2)$, $\operatorname{tr}(F \wedge F) = F_1 \wedge F_1 + F_2 \wedge F_2$ and $\operatorname{tr}F \wedge \operatorname{tr}F = (F_1 + F_2) \wedge (F_1 + F_2)$, so the difference is
$-2F_1 \wedge F_2$ (the $2$-forms commute) and $c_2 = -\frac{1}{4\pi^2}F_1 \wedge F_2 = \frac{i}{2\pi}F_1 \wedge \frac{i}{2\pi}F_2$, as predicted by the Whitney formula.
\end{solution}

\begin{exercise}\label{characteristic-classes:exercise-indices-plane}
Using Definition~\ref{characteristic-classes:definition-index} with the standard frame of $\RR^2 = \CC$, show that the vector field $X(z) = z^k$ (that is,
$X = \operatorname{Re}(z^k)\partial_x + \operatorname{Im}(z^k)\partial_y$) has index $k$ at $0$ for $k \geq 1$, and that $X(z) = \bar z^k$ has index $-k$. Which of these are gradients of functions?
\end{exercise}

\begin{solution}
On $S^1$, $\sigma = y^1dy^2 - y^2dy^1$, and on the circle $z = \varepsilon e^{i\varphi}$, $0 \leq \varphi \leq 2\pi$, positively oriented as the boundary of the disc, $z^k/|z^k| = e^{ik\varphi}$,
so $f(\varphi) = (\cos k\varphi, \sin k\varphi)$ and $f^*\sigma = k(\cos^2k\varphi + \sin^2k\varphi)\,d\varphi = k\,d\varphi$, with integral $2\pi k$; dividing by $|S^1| = 2\pi$ gives $k$. For $\bar z^k$,
$f(\varphi) = (\cos k\varphi, -\sin k\varphi)$ and $f^*\sigma = -k\,d\varphi$, giving $-k$. A gradient field $\nabla F = F_x\partial_x + F_y\partial_y$ corresponds to the complex function $F_x + iF_y$.
For $k = 1$, $\bar z = x - iy$ is the gradient of $\frac12(x^2 - y^2)$ (a saddle), and $z$ that of $\frac12(x^2 + y^2)$. In general $\bar z^k$ is the gradient of $\frac{1}{k+1}\operatorname{Re}(z^{k+1})$,
since for a holomorphic $h = u + iv$ the gradient of $u$ is $u_x + iu_y = u_x - iv_x = \overline{h^{\prime}}$; for $k = 2$ this is the monkey saddle
$\frac13\operatorname{Re}(z^3) = \frac13(x^3 - 3xy^2)$. On the other hand $z^k$ for
$k \geq 2$ is not a gradient: $F_x = \operatorname{Re}z^k$ and $F_y = \operatorname{Im}z^k$ would require $\partial_y\operatorname{Re}z^k = \partial_x\operatorname{Im}z^k$, but by the
Cauchy--Riemann equations $\partial_y\operatorname{Re}z^k = -\partial_x\operatorname{Im}z^k$, so both vanish, which fails for $k \geq 2$ (for $k = 2$, $\operatorname{Im}z^2 = 2xy$ has
$\partial_x = 2y \neq 0$).
\end{solution}

\begin{exercise}\label{characteristic-classes:exercise-linear-index}
Let $A \in \GL_n(\RR)$ and let $s(u) = Au$ on $\RR^n$, with the standard frame. Show that $\operatorname{ind}_0(s) = \operatorname{sgn}\det A$. (Write $A = hP$ with $h$ orthogonal and $P$
symmetric positive definite.)
\end{exercise}

\begin{solution}
$A^TA$ is symmetric positive definite, so by the spectral theorem it has a symmetric positive definite square root $P$, and $h = AP^{-1}$ satisfies $h^Th = P^{-1}A^TAP^{-1} = I$. The matrices
$P_t = (1 - t)I + tP$ are positive definite, so $hP_tu \neq 0$ for $u \neq 0$, and $f_t(u) = hP_tu/|hP_tu|$ is a homotopy on $S_\varepsilon$ from $hu/|u|$ to $Au/|Au|$. As in the proof of
Lemma~\ref{characteristic-classes:lemma-index}, homotopic maps give the same integral, so by Lemma~\ref{characteristic-classes:lemma-index}(2) the index is $\det h = \det A/\det P$, which has the
sign of $\det A$.
\end{solution}

\begin{exercise}\label{characteristic-classes:exercise-rotation-gauge}
Let $(e^{\prime}_1, e^{\prime}_2)$ be an oriented orthonormal frame of an oriented rank $2$ bundle with a metric connection, $\psi$ a smooth function, and
$e_1 = \cos\psi\,e^{\prime}_1 + \sin\psi\,e^{\prime}_2$, $e_2 = -\sin\psi\,e^{\prime}_1 + \cos\psi\,e^{\prime}_2$. Show that $\theta_{21} = \theta^{\prime}_{21} + d\psi$, and deduce again that $F_{12}$ does not
depend on the frame. Compare with the transformation rule $\theta = g^{-1}\theta^{\prime}g + g^{-1}dg$.
\end{exercise}

\begin{solution}
$\nabla e_1 = d\psi \otimes (-\sin\psi\,e^{\prime}_1 + \cos\psi\,e^{\prime}_2) + \cos\psi\,\nabla e^{\prime}_1 + \sin\psi\,\nabla e^{\prime}_2$, and $\nabla e^{\prime}_1 = \theta^{\prime}_{21} \otimes e^{\prime}_2$,
$\nabla e^{\prime}_2 = -\theta^{\prime}_{21} \otimes e^{\prime}_1$. So $\nabla e_1 = (d\psi + \theta^{\prime}_{21}) \otimes e_2$, that is $\theta_{21} = \theta^{\prime}_{21} + d\psi$. Then
$F_{12} = d\theta_{12} = -d\theta_{21} = -d\theta^{\prime}_{21} = F^{\prime}_{12}$. In matrix form $g$ is the rotation by $\psi$, $g^{-1}\theta^{\prime}g = \theta^{\prime}$ because rotations of the plane commute, and
$g^{-1}dg$ is the rotation generator $\bigl(\begin{smallmatrix} 0 & -1\\ 1 & 0\end{smallmatrix}\bigr)$ times $d\psi$, whose $(2,1)$ entry is $d\psi$.
\end{solution}

\begin{exercise}\label{characteristic-classes:exercise-one-zero-sphere}
On $S^2 = \CC\PP^1$ with the charts $z$ and $w = 1/z$, let $X$ be the vector field equal to $\partial_x$ in the chart $z = x + iy$. Show that $X$ extends smoothly to the point $w = 0$, where it
vanishes, and has index $2$ there. Deduce from Corollary~\ref{characteristic-classes:corollary-gauss-bonnet-chern}(1) that $\int_{S^2}e(TS^2, \nabla) = 2$ for every metric and every metric
connection, without using the Poincar\'e--Hopf theorem.
\end{exercise}

\begin{solution}
In complex notation $X$ is the constant field $1$ in the chart $z$: its flow is $z \mapsto z + t$. In the chart $w = 1/z$ this flow is $w \mapsto w/(1 + tw)$, whose derivative at $t = 0$ is
$-w^2$. So in the chart $w$ the field is $X(w) = -w^2$, which is smooth, also at $w = 0$, where it vanishes. The chart $w$ is positively oriented, because $z \mapsto 1/z$ is holomorphic, hence
orientation-preserving (its Jacobian determinant is $|{-1/z^2}|^2 > 0$). On the circle $w = \varepsilon e^{i\varphi}$ the direction of $-w^2$ is $-e^{2i\varphi} = e^{i(2\varphi + \pi)}$, which turns twice: the
index is $2$, by the computation of Exercise~\ref{characteristic-classes:exercise-indices-plane}, using Lemma~\ref{characteristic-classes:lemma-index}(1) to work with the coordinate frame. As $X$ has no
other zeros, Corollary~\ref{characteristic-classes:corollary-gauss-bonnet-chern}(1) gives $\int_{S^2}e(TS^2, \nabla) = 2$.
\end{solution}

\begin{exercise}\label{characteristic-classes:exercise-euler-form-sum}
Let $E_1$, $E_2$ be oriented bundles of even ranks with metrics and metric connections $\nabla_1$, $\nabla_2$, and give $E_1 \oplus E_2$ the orthogonal sum metric, the direct sum connection
and the orientation of an oriented frame of $E_1$ followed by one of $E_2$. Show that $e(\nabla_1 \oplus \nabla_2) = e(\nabla_1) \wedge e(\nabla_2)$. For a compact oriented surface $S$ with a metric,
deduce that $\int_{S \times S}e(T(S \times S), \nabla) = \chi(S)^2$ for the product metric and the direct sum of the Levi-Civita connections.
\end{exercise}

\begin{solution}
In the frame of $E_1 \oplus E_2$ made of oriented orthonormal frames of $E_1$ and $E_2$, the connection form and the curvature are block diagonal, $F = F_1 \oplus F_2$. By
Exercise~\ref{multilinear-algebra:exercise-pfaffian-blocks}, applied to matrices of commuting even forms, $\mathrm{Pf}(F) = \mathrm{Pf}(F_1) \wedge \mathrm{Pf}(F_2)$, and the factors $(2\pi)^{-m}$
split accordingly. For $S \times S$, $T(S \times S) = p_1^*TS \oplus p_2^*TS$, and the pullback connections have Euler forms $p_i^*e(TS, \nabla)$ (Proposition~\ref{characteristic-classes:proposition-euler-form}(4)), so
$\int_{S \times S}p_1^*e \wedge p_2^*e = (\int_Se)^2$ by Fubini, and $\int_Se(TS, \nabla) = \chi(S)$ by Corollary~\ref{characteristic-classes:corollary-gauss-bonnet-chern}.
\end{solution}
